\documentclass[11pt,a4paper]{article}
\usepackage[margin=1in]{geometry}
\usepackage{pbalgebras}

\title{Perfectly Bounded Real Banach Algebras\\
\large An involution-free, square-property extension of Gelfand duality}
\author{S. Pottel\thanks{With computational and editorial assistance from Claude
(Anthropic). Numerical claims are reproducible from the \texttt{verify/} package
of the accompanying repository.}}
\date{\today}

\begin{document}
\maketitle

\begin{abstract}
We study a class of unital real Banach algebras defined by two elementary
conditions --- a \emph{square property} $\norm{a^2}=\norm{a}^2$ and a
\emph{spectral reality} condition --- each relaxed by a \emph{commutator budget}
that vanishes on the centre.  On commutative algebras the unrelaxed conditions
characterise $C(X,\RR)$ for $X$ compact Hausdorff, and commutativity itself is a
theorem rather than a hypothesis; the relaxed conditions admit genuinely
noncommutative objects such as $M_2(\RR)$, $\HH$ and $T_2(\RR)$ while continuing
to exclude nilpotent thickenings such as $\RR[\varepsilon]/(\varepsilon^2)$.  The
resulting category $\PB$ is explicitly \emph{not} aimed at $C^*$-algebras: it
organises noncommutative geometry by normality (an eigenvalue radius) rather
than by an involution (a singular-value radius), and the two theories meet only
on the real forms of $C^*$-algebras.

Our main structural results are: every object carries a canonical compact
Hausdorff \emph{base} extracted from its centre; the morphism law forces
morphisms to preserve centres, so that the base is \emph{functorial}; and the
commutative part is a \emph{coreflective} subcategory of $\PB$ equivalent to
$\CompHaus\op$, exhibiting every object as a bundle of ``noncommutative points''
over a classical base.  We prove that the constant $C=\tfrac14$ is a phase
transition --- below it the theory is rigid and semisimple, at it the matrix and
triangular examples enter simultaneously --- and we settle, in the affirmative,
an open question from earlier notes about equalizer cones, showing that the
scalars are \emph{not} the equalizer of $\id$ and $\Ad_{\mathrm{diag}(1,-1)}$ on
$M_2(\RR)$.  Throughout, each statement is tagged \statusproved,
\statusnum, \statusconj\ or \statusopen, and the numerical tags are backed by a
reproducible verification package.
\end{abstract}

\tableofcontents

\section{Introduction}
\label{sec:intro}

\subsection{The fork}

Gelfand duality identifies commutative unital $C^*$-algebras with compact
Hausdorff spaces, and noncommutative geometry proceeds by dropping
commutativity while keeping the involution.  The $C^*$-identity
$\norm{a}^2=\norm{a^*a}$ makes the norm a \emph{singular-value} radius: it is
computed from $a$ and $a^*$ together, and it is insensitive to how badly $a$
fails to be normal.

There is a second, less travelled road out of the commutative case.  On a
commutative $C^*$-algebra one equally has
\begin{equation}
  \norm{a^2}=\norm{a}^2,
  \label{eq:sq-intro}
\end{equation}
and \eqref{eq:sq-intro} makes the norm an \emph{eigenvalue} radius: it says
exactly that $\norm{a}$ equals the spectral radius $r(a)$.  On a $C^*$-algebra
the two coincide precisely on the normal elements.  Dropping commutativity while
keeping \eqref{eq:sq-intro} --- rather than keeping the involution --- is the
fork taken in these notes.

Taken literally the fork is a dead end: \eqref{eq:sq-intro} together with a
spectral reality condition forces commutativity outright
(Theorem~\ref{thm:commutativity-forced}), so the ``noncommutative'' theory has
no objects.  The proposal studied here is to relax both conditions by a
\emph{commutator budget}: a quantity $\gauge(a)$, homogeneous of degree $2$,
which vanishes exactly on the centre and therefore leaves the commutative theory
untouched, but which funds a controlled failure of \eqref{eq:sq-intro} off the
centre.  The resulting objects we call \emph{perfectly bounded algebras}, or
\emph{PB-algebras}.

\subsection{What the budget buys}

The relaxed axioms read
\[
  \norm{a}^2-\norm{a^2}\le C\,\gauge(a),
  \qquad
  \bigl\lVert\,\norm{a^2}1-a^2\,\bigr\rVert\le\norm{a^2}+C\,\gauge(a),
\]
and they are genuinely two-parameter: the first measures \emph{approximate
normality}, the second \emph{approximate spectral reality}.  With the natural
gauge $\gauge_1(a)=\norm{\ad_a}^2$ and the constant $C=\tfrac14$ the class
contains $C(X,\RR)$, the quaternions $\HH$, the matrix algebras $M_n(\RR)$ and
the upper-triangular algebras $T_n(\RR)$, and it excludes
$\RR[\varepsilon]/(\varepsilon^2)$ --- the budget is zero on a commutative
algebra, so a nilpotent thickening cannot pay for its square defect.  The
theory therefore admits noncommutative points while refusing infinitesimal ones.

\subsection{Results}

The notes that this document supersedes recorded a collection of axioms,
examples and half-settled categorical questions.  The present treatment
reorganises that material around three structural theorems and corrects two
errors.

\begin{itemize}
\item \textbf{The base is functorial} (\S\ref{sec:base}).  On the centre both
  axioms collapse to the unrelaxed ones, so $\Zc(A)\cong C(X_A,\RR)$ for a
  compact Hausdorff $X_A$: every object carries a canonical classical base,
  \emph{extracted} from the algebra rather than imposed on it.  The morphism law
  --- contractive homomorphisms that do not increase the budget --- turns out to
  force $f(\Zc(A))\subseteq\Zc(B)$ automatically
  (Lemma~\ref{lem:centre-preserved}), so $A\mapsto X_A$ is a functor
  $\PB\to\CompHaus\op$.  This is the precise sense in which the category is
  ``fibred over classical geometry''.

\item \textbf{The commutative part is coreflective}
  (Theorem~\ref{thm:coreflection}).  The centre functor $A\mapsto\Zc(A)$ is
  right adjoint to the inclusion $\CommPB\hookrightarrow\PB$, and
  $\CommPB\simeq\CompHaus\op$.  Consequently every PB-algebra is a
  $C(X_A,\RR)$-algebra and so the section algebra of an upper semicontinuous
  field over its own base (\S\ref{sec:bundle}): PB-algebras are bundles of
  pointless objects over compact Hausdorff spaces.

\item \textbf{$C=\tfrac14$ is a phase transition}
  (Theorem~\ref{thm:phase-transition}).  For $C<\tfrac14$ every object satisfies
  $r(a)\ge(1-4C)\norm{a}$, hence is semisimple and contains no nonzero nilpotent
  element; at $C=\tfrac14$ the examples $M_2(\RR)$ and $T_2(\RR)$ enter
  simultaneously, and no choice of constant separates them
  (Proposition~\ref{prop:no-constant-separates}).  The value $\tfrac14$ is thus
  not a normalisation artefact but the location of the only interesting
  threshold in the family.
\end{itemize}

Two corrections to the earlier notes are recorded where they arise.  The
identification of the monomorphisms of $\PB$ with the budget-preserving
embeddings is withdrawn: injective morphisms are already monic, and the
budget-preserving embeddings are the correct notion of \emph{subobject}, not of
monomorphism (Remark~\ref{rem:mono-correction}).  The claim that the failure of
the inclusion $\Delta\hookrightarrow M_2(\RR)$ to be legal proves incompleteness
was already retracted there; we go further and \emph{answer} the open question
it left behind.  Theorem~\ref{thm:t2-cone} exhibits an explicit legal cone
$T_2(\RR)\to M_2(\RR)$ with non-scalar diagonal image, so the scalars are not
the equalizer of $\id$ and $\Ad_{\operatorname{diag}(1,-1)}$.  Whether
\emph{some} object is remains open, and \S\ref{sec:categorical} reduces that
question to the representability of an explicit functor.

\subsection{Status discipline}

Every substantive assertion carries one of four tags.  \statusproved\ means a
complete proof is given or a precise citation is supplied.  \statusnum\ means
the statement is a finite-dimensional inequality that has been checked by the
optimisation code in \texttt{verify/}, to the tolerance reported in
\S\ref{sec:status}, but not proved.  \statusconj\ and \statusopen\ mean what they
say.  The ledger in \S\ref{sec:status} collects all of them in one table; it is
the honest summary of where the theory stands, and it is intended to shrink in
the \statusopen\ column over time.

\section{Preliminaries and conventions}
\label{sec:prelim}

\begin{convention}\label{conv:standing}
Throughout, an \emph{algebra} means a unital real Banach algebra $A$ whose norm
is submultiplicative with $\norm{1}=1$.  Homomorphisms are unital and
$\RR$-linear.  We write $\Zc(A)$ for the centre, $\rad(A)$ for the Jacobson
radical, $[a,z]=az-za$ and $\ad_a=[a,-]$.  We allow the \emph{degenerate
algebra} $0$, in which $1=0$; see Convention~\ref{conv:degenerate}.
\end{convention}

\subsection{Complexification and spectrum}

For a real algebra $A$ let $A_\CC=A\otimes_\RR\CC$ carry any of the usual
equivalent norms, e.g.\
$\norm{x+iy}=\sup_{\theta}\norm{x\cos\theta-y\sin\theta}$, which is
submultiplicative and restricts to the given norm on $A$.  The \emph{spectrum}
of $a\in A$ is
\[
  \spec(a)=\{\lambda\in\CC:\lambda1-a\ \text{is not invertible in}\ A_\CC\},
\]
a nonempty compact subset of $\CC$ stable under complex conjugation, and the
spectral radius is $r(a)=\max\{\abs{\lambda}:\lambda\in\spec(a)\}$.  Beurling's
formula $r(a)=\lim_n\norm{a^n}^{1/n}$ holds verbatim, since the norms on $A$ and
$A_\CC$ are equivalent within a factor of $2$ and the limit is insensitive to
that.  A \emph{character} of $A$ is a nonzero homomorphism $\chi:A\to\CC$ of
real algebras with $\chi(1)=1$; every character is automatically contractive,
and $\abs{\chi(a)}\le r(a)$.

\begin{remark}
Nothing forces a noncommutative $A$ to have any characters at all: $M_2(\RR)$
has none, since a character would kill all commutators and
$M_2(\RR)=\RR1\oplus[M_2(\RR),M_2(\RR)]$, while the resulting functional
$a\mapsto\tfrac12\tr(a)$ is not multiplicative.  Statements below that quantify
over characters are therefore vacuous on such algebras; this is a feature, not a
gap.  The geometric content is carried by the centre instead (\S\ref{sec:base}).
\end{remark}

\subsection{Real Gelfand theory}

We use the real-scalar Gelfand theory in the form developed by Kulkarni and
Limaye \cite{KulkarniLimaye}.  For commutative $A$ the character space
$\Omega(A)$ of $A_\CC$ is compact Hausdorff and carries the involution
$\chi\mapsto\bar\chi$, $\bar\chi(a)=\overline{\chi(a)}$; the Gelfand transform
\[
  \Gamma:A\longrightarrow C(\Omega(A),\CC),\qquad \Gamma(a)(\chi)=\chi(a)
\]
has image inside the \emph{real function algebra} of conjugation-equivariant
functions, $\norm{\Gamma(a)}_\infty=r(a)$, and $\ker\Gamma=\rad(A)$.  The
involution is trivial --- so that the image consists of real-valued functions
--- exactly when every character is real-valued.  We write $\CompHaus$ for the
category of compact Hausdorff spaces and continuous maps.

\begin{theorem}[Real Stone--Weierstrass]\label{thm:stone-weierstrass}
A closed unital subalgebra of $C(X,\RR)$, $X$ compact Hausdorff, which separates
the points of $X$ is all of $C(X,\RR)$.
\end{theorem}

\subsection{Two classical inputs}

We shall use the following two results as black boxes; both are standard, and
both are the reason the rigid core of \S\ref{sec:rigid} is as rigid as it is.

\begin{theorem}[Hirschfeld--\.Zelazko \cite{HirschfeldZelazko}]
\label{thm:hz}
Let $B$ be a complex Banach algebra whose spectral radius is submultiplicative,
$r(xy)\le r(x)r(y)$ for all $x,y$.  Then $B/\rad(B)$ is commutative.
\end{theorem}

\begin{theorem}[Stampfli \cite{Stampfli}]\label{thm:stampfli}
For $a\in B(H)$, $H$ a complex Hilbert space,
$\norm{\ad_a}=2\dist(a,\CC1)$.
\end{theorem}

Theorem~\ref{thm:stampfli} is stated over $\CC$; the real analogue we use ---
$\norm{\ad_a}=2\dist(a,\RR1)$ for $a\in M_n(\RR)$ acting on $\RR^n$ --- is
recorded as Proposition~\ref{prop:real-stampfli} and is, in the present
treatment, verified numerically rather than proved.

\section{The rigid commutative core}
\label{sec:rigid}

\begin{definition}[The rigid axioms]\label{def:rigid}
An algebra $A$ satisfies
\begin{align}
  \tag{SQ}\label{ax:sq}
  &\norm{a^2}=\norm{a}^2 &&\text{for all } a\in A, \\
  \tag{SR}\label{ax:sr}
  &\bigl\lVert\,\norm{a^2}1-a^2\,\bigr\rVert\le\norm{a^2} &&\text{for all } a\in A .
\end{align}
We call \eqref{ax:sq} the \emph{square property} and \eqref{ax:sr} the
\emph{spectral reality condition}.
\end{definition}

\begin{remark}[On the names]
\eqref{ax:sr} appears in the earlier notes as ``(C2)''.  The name used here
records what it does: on a commutative algebra it is exactly the statement that
no character takes a non-real value (Lemma~\ref{lem:sr-characters}).  For
$A=C(X,\RR)$ both axioms are immediate, since $a^2\ge0$ and hence
$0\le\norm{a^2}-a^2(x)\le\norm{a^2}$ pointwise.
\end{remark}

\begin{lemma}[The square property is a spectral radius condition]
\label{lem:sq-radius}
$A$ satisfies \eqref{ax:sq} if and only if $r(a)=\norm{a}$ for all $a\in A$.  In
that case $A$ has no nonzero quasinilpotent elements; in particular
$\rad(A)=0$ and $A$ is semisimple.
\end{lemma}

\begin{proof}
If \eqref{ax:sq} holds then $\norm{a^{2^n}}=\norm{a}^{2^n}$ by induction, so
$r(a)=\lim_n\norm{a^{2^n}}^{1/2^n}=\norm{a}$.  Conversely if $r=\norm{\cdot}$
then $\norm{a}^2=r(a)^2=r(a^2)=\norm{a^2}$, using $\spec(a^2)=\spec(a)^2$.  A
quasinilpotent has $r=0$, hence norm $0$; and $\rad(A)$ consists of
quasinilpotents.
\end{proof}

\begin{lemma}[What \eqref{ax:sr} says about characters]\label{lem:sr-characters}
Let $A$ satisfy \eqref{ax:sr} and let $\chi$ be a character of $A$.  Then
$\chi(a)\in\RR$ for every $a\in A$.
\end{lemma}

\begin{proof}
Write $\chi(a)=\alpha+i\beta$ with $\alpha,\beta\in\RR$ and put
$b=a-\alpha1\in A$, so that $\chi(b)=i\beta$ and $\chi(b^2)=-\beta^2$.  Applying
$\chi$ to $\norm{b^2}1-b^2$ and using $\abs{\chi(x)}\le\norm{x}$,
\[
  \bigl\lVert\,\norm{b^2}1-b^2\,\bigr\rVert
  \;\ge\;\bigl\lvert\,\norm{b^2}+\beta^2\,\bigr\rvert
  \;=\;\norm{b^2}+\beta^2 .
\]
By \eqref{ax:sr} applied to $b$ the left side is at most $\norm{b^2}$, so
$\beta^2\le0$ and $\beta=0$.
\end{proof}

\begin{remark}
The translation $a\rightsquigarrow a-\alpha1$ is the whole trick, and it is the
step that the earlier notes performed only in the special case $\chi(a)=i$.  It
matters that \eqref{ax:sr} is assumed for \emph{all} elements; the condition at
a single $a$ says much less.
\end{remark}

\begin{theorem}[Real Gelfand duality]\label{thm:gelfand}
For a commutative algebra $A$ the following are equivalent.
\begin{enumerate}
\item[(i)] $A$ satisfies \eqref{ax:sq} and \eqref{ax:sr}.
\item[(ii)] $A\cong C(X,\RR)$ isometrically and unitally, for some compact
  Hausdorff $X$.
\end{enumerate}
Moreover $X$ is then homeomorphic to the character space of $A$, and
$C(-,\RR)$ and $\Spec$ are quasi-inverse contravariant equivalences between
$\CompHaus$ and the category of such algebras.
\end{theorem}

\begin{proof}
(ii)$\Rightarrow$(i) is the computation in the preceding remark.

(i)$\Rightarrow$(ii).  By Lemma~\ref{lem:sq-radius}, $r=\norm{\cdot}$ and
$\rad(A)=0$, so the Gelfand transform $\Gamma$ is injective and isometric onto
its image.  By Lemma~\ref{lem:sr-characters} every character is real-valued, so
the conjugation on $\Omega(A)$ is trivial and $\Gamma(A)$ is a closed unital
subalgebra of $C(X,\RR)$, where $X=\Omega(A)$.  Characters separate the points
of $X$ by definition, so $\Gamma(A)=C(X,\RR)$ by
Theorem~\ref{thm:stone-weierstrass}.  For the last clause, the characters of
$C(X,\RR)$ are the point evaluations, and a unital homomorphism
$C(X,\RR)\to C(Y,\RR)$ is $\varphi^*$ for a unique continuous
$\varphi:Y\to X$.
\end{proof}

\begin{theorem}[Commutativity is a theorem, not a hypothesis]
\label{thm:commutativity-forced}
Let $A$ be an algebra, not assumed commutative, satisfying \eqref{ax:sq} and
\eqref{ax:sr}.  Then $A$ is commutative, hence $A\cong C(X,\RR)$.
\end{theorem}

\begin{proof}
By Lemma~\ref{lem:sq-radius}, $r=\norm{\cdot}$ on $A$; in particular $r$ is
both subadditive and submultiplicative on $A$, and $A$ has no nonzero
quasinilpotent.  Granting the real-scalar form of Theorem~\ref{thm:hz} --- that
a real Banach algebra with submultiplicative spectral radius has commutative
quotient modulo its radical, see \cite[Ch.~3]{KulkarniLimaye} --- we get
$[x,y]\in\rad(A)$ for all $x,y\in A$.  Elements of the radical are
quasinilpotent, so $\norm{[x,y]}=r([x,y])=0$ and $A$ is commutative.  Now apply
Theorem~\ref{thm:gelfand}.
\end{proof}

\begin{remark}[Status of Theorem~\ref{thm:commutativity-forced}]
\label{rem:hz-status}
The proof is complete except that it invokes a real-scalar version of
Theorem~\ref{thm:hz}.  Over $\CC$ the statement is classical and the deduction
is immediate: a complex Banach algebra with $\norm{a^2}=\norm{a}^2$ throughout
has $r=\norm{\cdot}$, hence submultiplicative $r$ and zero radical, hence is
commutative and is a uniform algebra.  Over $\RR$ one wants either the cited
real form, or a reduction to the complex case; the obstacle to the naive
reduction is that $r=\norm{\cdot}$ on $A$ does not by itself give
$r=\norm{\cdot}$ on $A_\CC$.  We flag this in the ledger
(\S\ref{sec:status}) as \statusproved\ modulo the citation, and record closing
the gap by hand as Problem~\ref{prob:real-hz}.
\end{remark}

\begin{remark}[The finite-dimensional shadow]\label{rem:fd-shadow}
It is worth seeing the mechanism on a semisimple real algebra
$A=\prod_i M_{n_i}(D_i)$ with $D_i\in\{\RR,\CC,\HH\}$.  The square property
kills every block with $n_i\ge2$, because such a block contains a nonzero
square-zero element, which \eqref{ax:sq} forces to vanish.  Spectral reality
kills the $\CC$ and $\HH$ directions, because in $\CC$ the element $i$ has
$i^2=-1$ and $\norm{\,\norm{i^2}1-i^2}=2>1=\norm{i^2}$.  What survives is
$\RR^j=C(\{1,\dots,j\},\RR)$.  Theorem~\ref{thm:commutativity-forced} is the
infinite-dimensional form of this stripping.
\end{remark}

\begin{remark}[Why a relaxation is needed at all]
Theorem~\ref{thm:commutativity-forced} is the obstruction that the rest of these
notes works around.  Any theory that keeps \eqref{ax:sq} and \eqref{ax:sr} on
the nose is the classical theory; a noncommutative extension must pay for
noncommutativity somewhere, and the proposal is to pay for it with a quantity
that is invisible on the centre.
\end{remark}

\section{Commutator gauges}
\label{sec:gauges}

\begin{definition}[Gauges]\label{def:gauges}
For $a\in A$ and $q\in[1,\infty]$ set
\[
  \gauge_1(a)=\sup_{\norm{z}\le1}\norm{[a,z]}^2=\norm{\ad_a}^2,
  \qquad
  \gauge_q(a)=\sup_{\norm{z}\le1}\norm{[a,z]^q}^{2/q},
  \qquad
  \gauge_\infty(a)=\sup_{\norm{z}\le1}r\bigl([a,z]\bigr)^2 .
\]
A \emph{gauge} means one of these.  The value $\gauge(a)$ is the
\emph{commutator budget} of $a$.
\end{definition}

\begin{lemma}[Elementary properties]\label{lem:gauge-basic}
Let $\gauge$ be any of the gauges.  For all $a\in A$ and $t,\lambda\in\RR$:
\begin{enumerate}
\item[(a)] \emph{(Homogeneity)} $\gauge(ta)=t^2\gauge(a)$.
\item[(b)] \emph{(Translation invariance)} $\gauge(a+\lambda1)=\gauge(a)$.
\item[(c)] \emph{(Universal bound)} $\gauge(a)\le4\norm{a}^2$, and more
  precisely $\gauge(a)\le4\dist(a,\Zc(A))^2$.
\item[(d)] \emph{(Ordering)} $\gauge_\infty(a)\le\gauge_q(a)\le\gauge_1(a)$ for
  every $q\in[1,\infty)$.
\item[(e)] \emph{(Vanishing locus)} $\gauge_1(a)=0$ if and only if
  $a\in\Zc(A)$.  For $q\ge2$ and for $q=\infty$ one has only the implication
  $a\in\Zc(A)\Rightarrow\gauge(a)=0$, and the converse fails.
\end{enumerate}
\end{lemma}

\begin{proof}
(a) and (b) are immediate from $\ad_{a+\lambda1}=\ad_a$ and bilinearity.
For (c), if $c\in\Zc(A)$ then $[a,z]=[a-c,z]$, so
$\norm{[a,z]}\le2\norm{a-c}\norm{z}$; submultiplicativity gives
$\norm{[a,z]^q}\le\norm{[a,z]}^q$, and raising to the power $2/q$ cancels the
exponent, so the bound $4\dist(a,\Zc(A))^2$ is uniform in $q$; for
$\gauge_\infty$ use $r\le\norm{\cdot}$.  (d) is $r(b)\le\norm{b^q}^{1/q}\le\norm{b}$.
For (e), $\gauge_1(a)=0$ says $\ad_a=0$, i.e.\ $a\in\Zc(A)$.  For finite
$q\ge2$, $\gauge_q(a)=0$ says only that $[a,z]^q=0$ for every $z$, and for
$q=\infty$ only that every $[a,z]$ is quasinilpotent; the element
$a=\operatorname{diag}(1,-1)\in T_2(\RR)$ is not central, yet every commutator
$[a,z]$ with $z\in T_2(\RR)$ is a multiple of $E_{12}$ and hence square-zero, so
$\gauge_q(a)=\gauge_\infty(a)=0$ for all $q\ge2$.
\end{proof}

\begin{remark}
The exponent $2$ is forced: both sides of the axioms of \S\ref{sec:axioms} must
be homogeneous of degree $2$ in $a$ for the resulting class to be closed under
rescaling of elements, and $\norm{a}^2-\norm{a^2}$ has degree $2$.
\end{remark}

\subsection{The gauge as a distance to the centre}

Part (c) of Lemma~\ref{lem:gauge-basic} bounds $\gauge_1$ by the distance to the
centre.  In the two basic noncommutative examples the bound is an equality, and
this is what makes the constants of \S\ref{sec:axioms} computable.

\begin{proposition}[Real Stampfli identity]\label{prop:real-stampfli}
For $A=M_n(\RR)$ with the operator norm and for $A=\HH$,
\[
  \gauge_1(a)=4\dist(a,\RR1)^2 .
\]
\statusnum\ for $M_n(\RR)$, $n\le4$, to a relative tolerance of $10^{-6}$;
\statusproved\ for $\HH$.  The complex statement for $B(H)$ is
Theorem~\ref{thm:stampfli}.
\end{proposition}

\begin{proof}[Proof for $\HH$]
Write $a=\alpha+v$ with $\alpha\in\RR$ and $v$ purely imaginary, so
$\dist(a,\RR1)=\abs{v}$.  Then $[a,z]=[v,z]$, and for $z$ a unit quaternion
orthogonal to $v$ one computes $[v,z]=2vz$, of modulus $2\abs{v}$; since
$\norm{[v,z]}\le2\abs{v}\norm{z}$ always, the supremum is $2\abs{v}$.
\end{proof}

\begin{proposition}[The identity fails on triangular algebras]
\label{prop:stampfli-fails-T}
For $A=T_n(\RR)$, $n\ge2$, the inequality $\gauge_1(a)\le4\dist(a,\RR1)^2$ is
strict for generic $a$.  \statusnum
\end{proposition}

\noindent
The reason is structural rather than numerical: $\gauge_1$ is a supremum over
$z$ \emph{in the algebra}, and $T_n(\RR)$ has a much smaller supply of unit
elements than $M_n(\RR)$ does.  This ambient-referencing character of the gauge
is the source of every difficulty in \S\ref{sec:categorical}: an equation can
delete the very elements $z$ that realise a budget.  We return to this in
Remark~\ref{rem:ambient-referencing}.

\begin{example}[The fundamental computation]\label{ex:E12}
Let $E_{12}\in M_2(\RR)$ be the matrix unit and $u=\operatorname{diag}(1,-1)$.
Then $\norm{E_{12}}=1$, $E_{12}^2=0$, and
$[E_{12},u]=-2E_{12}$, so
\[
  \gauge_1(E_{12})=4,\qquad \gauge_\infty(E_{12})=1,\qquad
  \gauge_q(E_{12})=1\ \ (q\ge2).
\]
The value $\gauge_1(E_{12})=4$ is the universal bound of
Lemma~\ref{lem:gauge-basic}(c), attained.  The drop at $q\ge2$ happens because
the budget-realising commutator $-2E_{12}$ is nilpotent and is therefore
annihilated by the $q$-th power; the supremum is then realised instead by
$[E_{12},E_{21}]=E_{11}-E_{22}$, of spectral radius $1$.  The same element
$E_{12}$, viewed inside $T_2(\RR)$, still has $\gauge_1=4$ --- the witness $u$
is upper triangular --- but $\gauge_\infty(E_{12})=0$ there, since every
commutator of upper-triangular matrices is strictly upper triangular, hence
nilpotent.  All four values are confirmed numerically in
\texttt{verify/tests/test\_gauges.py}.
\end{example}

\begin{remark}[$\gauge_1$ versus $\gauge_\infty$: discrimination against closure]
\label{rem:gauge-tradeoff}
Example~\ref{ex:E12} already contains the central trade-off of the whole
construction.  The spectral gauge $\gauge_\infty$ is \emph{discriminating}: it
sees only non-quasinilpotent commutators, so it distinguishes $M_2(\RR)$ from
$T_2(\RR)$ and reproduces, on the nose, the $C^*$-style exclusion of the
radical.  But it vanishes on whole directions, and quotients can manufacture
such vanishing, so a $\gauge_\infty$-class has no reason to be closed under
quotients (Proposition~\ref{prop:dinf-fragile}).  The norm gauge $\gauge_1$ is
\emph{closure-friendly}: it vanishes only on the centre, which is exactly what
makes the base functorial (\S\ref{sec:base}).

The earlier notes add that $\gauge_1$ pays for this by being unable to
distinguish $T_2(\RR)$ from $M_2(\RR)$ at any constant, and present the choice
as a forced trade between discrimination and closure.  That is not so:
$\kap_{\gauge_1}(T_2(\RR))=\kap^*>\tfrac14=\kap_{\gauge_1}(M_2(\RR))$
(Proposition~\ref{prop:T2-constant}), so constants in $[\tfrac14,\kap^*)$
discriminate exactly as $\gauge_\infty$ does on this example, while keeping the
functorial base.  How far that goes is
Conjecture~\ref{conj:semisimple}.  We develop the theory for $\gauge_1$, both
because functoriality of the base is the property that makes $\PB$ a category
worth having, and because the discrimination it was thought to lack is at least
partly available after all.
\end{remark}

\begin{warning}\label{warn:vanishing}
The earlier notes assert that every $\gauge_q$ ``vanishes exactly on directions
commuting with all of $A$''.  Lemma~\ref{lem:gauge-basic}(e) shows this is true
only for $q=1$: already $\operatorname{diag}(1,-1)\in T_2(\RR)$ is a
non-central element with $\gauge_q=0$ for every $q\ge2$.  The consequences are
not cosmetic.  It is precisely the equality $\{\gauge_1=0\}=\Zc(A)$ that makes
morphisms preserve centres (Lemma~\ref{lem:centre-preserved}) and hence makes
the classical base a functor; for $q\ge2$ that argument breaks, and with it the
geometric content of \S\ref{sec:base}.
\end{warning}

\section{The axioms}
\label{sec:axioms}

\begin{definition}[PB-algebras]\label{def:pb}
Fix a gauge $\gauge$ and a constant $C>0$.  An algebra $A$ is a
\emph{perfectly bounded algebra}, or \emph{PB-algebra}, for the data
$(\gauge,C)$ if for all $a\in A$
\begin{align}
  \tag{PB1}\label{ax:pb1}
  \norm{a}^2-\norm{a^2}&\le C\,\gauge(a), \\
  \tag{PB2}\label{ax:pb2}
  \bigl\lVert\,\norm{a^2}1-a^2\,\bigr\rVert&\le\norm{a^2}+C\,\gauge(a) .
\end{align}
We write $\PB^C_\gauge$ for the resulting category (morphisms in
\S\ref{sec:morphisms}), abbreviated $\PB$ when the data are understood, and
$\PB_\gauge=\bigcup_{C>0}\PB^C_\gauge$.  The unrelaxed case, where the budget
term is deleted from \eqref{ax:pb1} only, is written $\PB^0$.
\end{definition}

\begin{convention}[The degenerate object]\label{conv:degenerate}
We admit the zero algebra $0$, in which $1=0$, as an object of $\PB^C_\gauge$;
both axioms hold vacuously.  This is the same convention that makes the
category of unital $C^*$-algebras complete, and without it $\PB$ has no
terminal object and fails to have limits for a reason that has nothing to do
with the subject (Proposition~\ref{prop:terminal}).
\end{convention}

\subsection{The axioms as an invariant}

Definition~\ref{def:pb} quantifies over a constant, which is awkward: the
constant is not an invariant of $A$, only a bound on one.  The following
repackaging is equivalent and is what the verification code computes.

\begin{definition}[Critical constants]\label{def:critical}
For an algebra $A$ and a gauge $\gauge$ put
\[
  \kap_\gauge(A)=\sup_{a\in A}\frac{\norm{a}^2-\norm{a^2}}{\gauge(a)},
  \qquad
  \lam_\gauge(A)=\sup_{a\in A}\frac{\bigl(\bigl\lVert\norm{a^2}1-a^2\bigr\rVert-\norm{a^2}\bigr)_+}{\gauge(a)},
\]
with the conventions $0/0=0$ and $t/0=+\infty$ for $t>0$.  Set
$C_\gauge(A)=\max\bigl(\kap_\gauge(A),\lam_\gauge(A)\bigr)$.
\end{definition}

\begin{lemma}\label{lem:critical-reformulation}
$A\in\PB^C_\gauge$ if and only if $C_\gauge(A)\le C$.  In particular
$A\in\PB_\gauge$ if and only if $C_\gauge(A)<\infty$, and this requires in
particular that every $a$ with $\gauge(a)=0$ satisfy \eqref{ax:sq} and
\eqref{ax:sr} exactly.
\end{lemma}

\begin{proof}
Immediate from the definitions; both defects and the gauge are homogeneous of
degree $2$, so the suprema may equivalently be taken over the unit sphere.
\end{proof}

\begin{remark}[Why this matters]
The reformulation makes three things visible that Definition~\ref{def:pb}
obscures.  First, $C_\gauge(-)$ is a genuine numerical invariant of an algebra,
computable in finite dimensions, and the categories $\PB^C_\gauge$ are a
filtration of $\PB_\gauge$ by its value.  Second, the case $\gauge(a)=0$ is not
a boundary case to be waved through but the heart of the matter: on the
$\gauge_1$-null elements --- that is, on the centre --- the axioms are the
rigid ones, with no slack whatsoever.  Third, it is now a well-posed finite
computation to ask whether a given finite-dimensional algebra is a PB-algebra,
which is what \texttt{verify/} does.
\end{remark}

\subsection{The two defects}

\begin{remark}[Reading the axioms]\label{rem:reading}
\eqref{ax:pb1} measures \emph{approximate normality} and \eqref{ax:pb2}
\emph{approximate spectral reality}, each with slack proportional to the
commutation defect.  The contrast with the $C^*$-identity is exact.  On a
$C^*$-algebra,
\[
  \norm{a}^2=\norm{a^*a}=r(a^*a)
  \qquad\text{whereas}\qquad
  \norm{a}^2=\norm{a^2}\iff a\ \text{normal},
\]
so the $C^*$-norm is a singular-value radius and the PB-norm an eigenvalue
radius.  The two agree exactly on normal elements; in the commutative case
everything is normal and both roads lead to $\CompHaus$.  ``Singular value
versus eigenvalue'' is the fork, and it is the reason the two theories are
transverse rather than nested (\S\ref{sec:cstar}).
\end{remark}

\begin{convention}[Normalisation]\label{conv:normalisation}
By Lemma~\ref{lem:gauge-basic}(c) one may prefer the normalised gauge
$\tilde\gauge=\tfrac14\gauge$, for which $\tilde\gauge(a)\le\norm{a}^2$ and
both sides of \eqref{ax:pb1} live on the common scale $[0,\norm{a}^2]$.  Then
$C=\tfrac14$ becomes $C=1$.  We keep the unnormalised gauge and carry the
$\tfrac14$, because the numerical value $\tfrac14$ is where the interesting
threshold sits (Theorem~\ref{thm:phase-transition}) and it is clearer to have
the threshold in the constant than hidden in the gauge.
\end{convention}

\section{First consequences}
\label{sec:consequences}

Throughout this section $A\in\PB^C_{\gauge_1}$.

\subsection{The centre is classical}

\begin{theorem}[The classical base]\label{thm:base}
The centre $\Zc(A)$ is a closed commutative subalgebra on which \eqref{ax:sq}
and \eqref{ax:sr} hold exactly.  Consequently
\[
  \Zc(A)\;\cong\;C(X_A,\RR)
\]
isometrically and unitally, for a compact Hausdorff space
$X_A=\Spec\Zc(A)$, uniquely determined up to homeomorphism.
\end{theorem}

\begin{proof}
The centre of a Banach algebra is a closed subalgebra containing $1$, and it is
commutative.  For $a\in\Zc(A)$ we have $\gauge_1(a)=0$ by
Lemma~\ref{lem:gauge-basic}(e), so \eqref{ax:pb1} reads
$\norm{a}^2\le\norm{a^2}$, which together with submultiplicativity gives
\eqref{ax:sq}; and \eqref{ax:pb2} reads exactly \eqref{ax:sr}.  Both are
statements about norms of elements of $A$, hence hold in $\Zc(A)$ with the
induced norm.  Now apply Theorem~\ref{thm:gelfand}.
\end{proof}

\begin{definition}\label{def:base}
$X_A$ is the \emph{classical base} of $A$.  An object with $X_A$ a single
point, equivalently $\Zc(A)=\RR1$, is called \emph{pointless}.
\end{definition}

\begin{remark}[Extracted, not imposed]\label{rem:extracted}
The base is read off the algebra; one is not required to sit over an ambient
space chosen in advance.  For $A=M_2(\RR)$ or $A=\HH$ the base is a single
point, so genuinely pointless noncommutative objects are admitted --- the
theory does not secretly reduce to bundles over a nontrivial space.  For
$A=C(X,\RR)$ the base is $X$ and the algebra is its own centre.  The content of
\S\ref{sec:base} is that this assignment is functorial, and \S\ref{sec:bundle}
that it exhibits every object as a bundle of pointless ones.
\end{remark}

\subsection{Approximate spectral reality}

\begin{theorem}[Characters are approximately real]\label{thm:approx-real}
Let $\chi$ be a character of $A$.  Then for every $a\in A$,
\[
  \bigl(\Ima\chi(a)\bigr)^2\;\le\;C\,\gauge_1(a).
\]
In particular every character of $A$ is real-valued on $\Zc(A)$, and if $A$ is
commutative then every character is real-valued.
\end{theorem}

\begin{proof}
Write $\chi(a)=\alpha+i\beta$ and set $b=a-\alpha1$, so $\chi(b^2)=-\beta^2$ and
$\gauge_1(b)=\gauge_1(a)$ by Lemma~\ref{lem:gauge-basic}(b).  As in
Lemma~\ref{lem:sr-characters},
$\bigl\lVert\norm{b^2}1-b^2\bigr\rVert\ge\norm{b^2}+\beta^2$, and
\eqref{ax:pb2} applied to $b$ bounds the left side by
$\norm{b^2}+C\gauge_1(b)$.
\end{proof}

\begin{remark}
This is the quantitative form of ``the theory has real spectra''.  Off the
centre the bound degrades exactly in proportion to the budget, which is what
lets $\HH$ and the rotation matrix $J=\begin{psmallmatrix}0&-1\\1&0\end{psmallmatrix}$
into the theory: both have $J^2=-1$, an unambiguously non-real spectral
direction, and both are funded by $\gauge_1(J)=4$.  The commutative algebra
$\CC$, which has the same spectral defect and \emph{no} budget, is excluded.
The slogan is that complex directions are allowed only when they are
noncommutative ones.
\end{remark}

\subsection{The phase transition at $C=\tfrac14$}

\begin{theorem}[Quantitative normaloid bound]\label{thm:normaloid}
Let $A\in\PB^C_{\gauge_1}$ with $C<\tfrac14$.  Then for every $a\in A$,
\[
  r(a)\;\ge\;(1-4C)\,\norm{a}.
\]
\end{theorem}

\begin{proof}
By \eqref{ax:pb1} and the universal bound $\gauge_1(x)\le4\norm{x}^2$,
\[
  \norm{x^2}\;\ge\;\norm{x}^2-C\gauge_1(x)\;\ge\;(1-4C)\norm{x}^2
  \qquad\text{for every } x\in A .
\]
Apply this with $x=a^{2^n}$ and put $u_n=\norm{a^{2^n}}^{1/2^n}$, so that
$u_{n+1}^{2^{n+1}}\ge(1-4C)\,u_n^{2^{n+1}}$, i.e.\
$u_{n+1}\ge(1-4C)^{1/2^{n+1}}u_n$.  Iterating from $u_0=\norm{a}$,
\[
  u_n\;\ge\;\norm{a}\prod_{k=1}^{n}(1-4C)^{2^{-k}}
  \;=\;\norm{a}\,(1-4C)^{\,1-2^{-n}}\;\ge\;\norm{a}\,(1-4C),
\]
and $r(a)=\lim_n u_n$ by Beurling's formula.
\end{proof}

\begin{corollary}\label{cor:subcritical-rigid}
If $C<\tfrac14$ then every $A\in\PB^C_{\gauge_1}$ is semisimple and contains no
nonzero nilpotent and no nonzero quasinilpotent element.
\end{corollary}

\begin{proof}
A quasinilpotent has $r(a)=0$, hence $\norm{a}=0$ by
Theorem~\ref{thm:normaloid}; nilpotents are quasinilpotent, and $\rad(A)$
consists of quasinilpotents.
\end{proof}

\begin{theorem}[Phase transition]\label{thm:phase-transition}
For the gauge $\gauge_1$:
\begin{enumerate}
\item[(a)] For every $C<\tfrac14$, the class $\PB^C_{\gauge_1}$ consists of
  semisimple algebras with $r\ge(1-4C)\norm{\cdot}$, and contains no algebra
  with a nonzero nilpotent --- in particular neither $M_2(\RR)$ nor $T_2(\RR)$.
\item[(b)] At $C=\tfrac14$ the pointless objects $M_2(\RR)$ and $\HH$ enter,
  with $\kap_{\gauge_1}(M_2(\RR))=\lam_{\gauge_1}(M_2(\RR))=\tfrac14$ attained
  at $E_{12}$ and at $J$ respectively.
\end{enumerate}
Part (a) is \statusproved; the equalities in (b) are \statusnum, the
inequalities $\ge\tfrac14$ being Examples~\ref{ex:H} and \ref{ex:M2}.
\end{theorem}

\begin{proof}[Proof of (a)]
The first clause is Theorem~\ref{thm:normaloid} and
Corollary~\ref{cor:subcritical-rigid}.  Both $M_2(\RR)$ and $T_2(\RR)$ contain
the nonzero nilpotent $E_{12}$, so neither is subcritical.
\end{proof}

\begin{theorem}[A constant \emph{does} separate $T_2$ from $M_2$]
\label{thm:constant-separates}
$\kap_{\gauge_1}(T_2(\RR))\ge\kap^*=0.3098421747\ldots>\tfrac14$, where
$\kap^*$ is the algebraic number of Proposition~\ref{prop:T2-constant}.  Hence
$T_2(\RR)\notin\PB^{1/4}_{\gauge_1}$, while $M_2(\RR)\in\PB^{1/4}_{\gauge_1}$,
and every constant $C\in[\tfrac14,\kap^*)$ admits the one and excludes the
other.
\end{theorem}

\begin{proof}
Proposition~\ref{prop:T2-constant} and Example~\ref{ex:M2}.
\end{proof}

\begin{remark}[Correction to the earlier notes]\label{rem:separation-correction}
Theorem~\ref{thm:constant-separates} contradicts the assertion of the earlier
notes that no $\gauge_1$-constant separates $T_2(\RR)$ from $M_2(\RR)$, an
assertion that rested on locating the extremiser of $T_2(\RR)$ at the shared
element $E_{12}$.  It does not sit there; see
Remark~\ref{rem:T2-correction}.  The consequences are structural, not
bookkeeping.  The earlier notes present the choice between $\gauge_1$ and
$\gauge_\infty$ as a forced trade --- $\gauge_1$ is closure-friendly and makes
the base functorial but cannot exclude non-central radical, $\gauge_\infty$
excludes it but vanishes on whole directions and so breaks closure.  At
$C=\tfrac14$ the trade is not forced in the one case we can compute: the norm
gauge excludes $T_2(\RR)$ \emph{and} keeps the functorial base of
\S\ref{sec:base}.  How general that is, we do not know, and it is the most
interesting question the correction opens up.
\end{remark}

\begin{conjecture}[Semisimplicity at the critical constant]
\label{conj:semisimple}
Every $A\in\PB^{1/4}_{\gauge_1}$ is semisimple, i.e.\ $\rad(A)=0$.
\end{conjecture}

\begin{remark}[Evidence]\label{rem:semisimple-evidence}
Note first that Conjecture~\ref{conj:semisimple} does \emph{not} follow from
Theorem~\ref{thm:normaloid}, which is vacuous at $C=\tfrac14$, and that it is
not a statement about nilpotent \emph{elements}: $M_2(\RR)\in\PB^{1/4}$ contains
the nilpotent $E_{12}$ and is nonetheless simple.

The evidence is that every elementary radical-carrying algebra we can compute is
excluded at $\tfrac14$.  For $\RR[\varepsilon]/(\varepsilon^2)$ the defect is
unfunded outright (Example~\ref{ex:excluded}).  For the others it is a
quantitative matter, and to exclude an algebra it suffices to exhibit one
element with ratio above $\tfrac14$ --- no supremum is needed.  Such witnesses
are found by \texttt{verify/experiments/radical\_scan.py}:
\begin{center}
\begin{tabular}{llc}
\toprule
Algebra & witness & $\bigl(\norm{a}^2-\norm{a^2}\bigr)/\gauge_1(a)$ \\
\midrule
$T_2(\RR)$   & $x_{0.45}$ & $0.30894$ \\
$T_3(\RR)$   & $x_{0.45}$ in the $\{1,3\}$ corner & $0.30894$ \\
$P_{1,2}(\RR)$ & $x_{0.45}$ in the $\{1,2\}$ corner & $0.30190$ \\
$P_{2,1}(\RR)$ & $x_{0.45}$ in the $\{1,3\}$ corner & $0.30047$ \\
\bottomrule
\end{tabular}
\end{center}
Here $x_c=\begin{psmallmatrix}c&1-c^2\\0&-c\end{psmallmatrix}$ is the traceless
family of Proposition~\ref{prop:T2-constant} and $P_{1,2}(\RR)$, $P_{2,1}(\RR)$
are the block upper-triangular (parabolic) subalgebras of $M_3(\RR)$.  The
parabolic cases are the informative ones: unlike $T_n(\RR)$ they contain a full
$M_2(\RR)$ block on the diagonal, so they have many more test elements with
which to fund a radical defect --- and they still cannot fund it at $\tfrac14$.
All four entries are \statusnum; each is stable to $10^{-4}$ under a tripling of
the optimisation effort, and the first is corroborated by the exact computation
of Proposition~\ref{prop:T2-constant}.

If the conjecture holds, then at the critical constant the norm gauge achieves
without an involution what the $C^*$-identity achieves with one --- the
exclusion of the radical --- while retaining the functoriality that
$\gauge_\infty$ and every $\gauge_q$, $q\ge2$, destroy
(Warning~\ref{warn:vanishing}), and the ``discrimination versus closure''
trade-off of Remark~\ref{rem:gauge-tradeoff} dissolves at exactly one value of
$C$.  This is Problem~\ref{prob:semisimple}.
\end{remark}

\begin{remark}[What the transition means]
The interval $C\in(0,\tfrac14)$ is not a family of interesting theories: by
Corollary~\ref{cor:subcritical-rigid} every object there is semisimple and
nilpotent-free, and the surviving finite-dimensional objects are products of
copies of $\RR$ and $\HH$ (Proposition~\ref{prop:strict-survivors}).  The
noncommutative content appears at $C=\tfrac14$, where the pointless objects
$M_2(\RR)$ and $\HH$ arrive together.  Above $\kap^*$ the non-central radical
arrives too.  So the constant is not a free parameter to be tuned away: it
indexes a genuine filtration, with at least two thresholds,
$\tfrac14$ and $\kap^*$, and the interval between them is where the theory looks
most like a noncommutative geometry --- noncommutative points, no radical, and a
functorial classical base.
\end{remark}

\begin{proposition}[The survivors at zero budget]\label{prop:strict-survivors}
Let $A$ be finite-dimensional and satisfy \eqref{ax:sq} exactly together with
\eqref{ax:pb2}.  Then $A$ is a product of copies of $\RR$ and $\HH$.
\end{proposition}

\begin{proof}
\eqref{ax:sq} gives $\rad(A)=0$ (Lemma~\ref{lem:sq-radius}), so by
Wedderburn $A\cong\prod_iM_{n_i}(D_i)$ with $D_i\in\{\RR,\CC,\HH\}$.  A block
with $n_i\ge2$ contains a nonzero square-zero element, contradicting
\eqref{ax:sq}.  A block $\CC$ is central, so its budget vanishes and
\eqref{ax:pb2} reduces to \eqref{ax:sr}, which $i\in\CC$ violates as in
Remark~\ref{rem:fd-shadow}.  Blocks $\RR$ and $\HH$ both survive: the norm of
$\HH$ is multiplicative, giving \eqref{ax:sq}, and
$\gauge_1(v)=4\abs{v}^2$ for $v$ purely imaginary
(Proposition~\ref{prop:real-stampfli}) funds \eqref{ax:pb2} at $C=\tfrac14$
with equality.
\end{proof}

\begin{remark}[Dyson's threefold way, with the unitary class excised]
Proposition~\ref{prop:strict-survivors} says the zero-budget theory sees the
orthogonal and symplectic classes and not the unitary one.  The
infinite-dimensional analogue --- that \eqref{ax:sq} together with
\eqref{ax:pb2} forces an isometric embedding into an algebra of sections with
$\RR$- and $\HH$-fibres --- would be a quaternionic Stone--Weierstrass theorem.
It is \statusconj; see Problem~\ref{prob:quaternionic-sw}.
\end{remark}

\section{Morphisms}
\label{sec:morphisms}

The gauge is not functorial: an arbitrary unital homomorphism can increase a
commutator budget, as the diagonal embedding $\RR^2\hookrightarrow M_2(\RR)$
already shows.  Rather than clamping the objects, one restricts the arrows.

\begin{definition}[PB-morphisms]\label{def:pb-morphism}
A \emph{PB-morphism} $f:A\to B$ is a unital algebra homomorphism which is
\begin{enumerate}
\item[(i)] contractive: $\norm{f(a)}\le\norm{a}$ for all $a$; and
\item[(ii)] budget-contracting: $\gauge_B(f(a))\le\gauge_A(a)$ for all $a$.
\end{enumerate}
\end{definition}

\begin{remark}[Terminology]
These are the arrows written $\dagger$ in the earlier notes.  We avoid that
symbol: it is standard for dagger categories, where it denotes an involutive
contravariant identity-on-objects functor, and the present notion is neither
involutive nor contravariant --- and the whole point of the theory is that
there is no involution.
\end{remark}

\begin{proposition}\label{prop:category}
PB-morphisms contain the identities and are closed under composition, so
$\PB^C_\gauge$ is a category.  Every quotient map $A\to A/I$ by a closed
two-sided ideal is budget-contracting and contractive.
\end{proposition}

\begin{proof}
Identities and composites are clear from the two inequalities.  For a quotient
$\pi:A\to A/I$, contractivity is the definition of the quotient norm.  For the
budget, let $\bar z\in A/I$ with $\norm{\bar z}<1$ and choose a lift $z$ with
$\norm{z}<1$; then
$\norm{[\pi(a),\bar z]}=\norm{\pi([a,z])}\le\norm{[a,z]}$, and similarly for
the $q$-th powers and for spectral radii, using that $\pi$ is a contractive
homomorphism.  Taking suprema gives
$\gauge_{A/I}(\pi(a))\le\gauge_A(a)$.
\end{proof}

\subsection{Isometries preserve the budget}

The following simple lemma replaces a claim of the earlier notes and is used
repeatedly.

\begin{lemma}[Isometric implies budget-nondecreasing]\label{lem:isometric}
Let $f:A\to B$ be an isometric unital homomorphism.  Then
$\gauge_B(f(a))\ge\gauge_A(a)$ for every $a\in A$.  Consequently an isometric
PB-morphism preserves the gauge exactly:
$\gauge_B\circ f=\gauge_A$.
\end{lemma}

\begin{proof}
For $z\in A$ with $\norm{z}\le1$ we have $\norm{f(z)}=\norm{z}\le1$ and
$f([a,z])=[f(a),f(z)]$, so
$\norm{[f(a),f(z)]}=\norm{[a,z]}$; the supremum defining $\gauge_B(f(a))$ is
over a larger set of test elements than the one defining $\gauge_A(a)$.  The
same argument applies verbatim to $\gauge_q$ and $\gauge_\infty$, using that an
isometric homomorphism preserves norms of powers and spectral radii of elements
in its image.  The final clause combines this with
Definition~\ref{def:pb-morphism}(ii).
\end{proof}

\begin{definition}[Saturated embeddings]\label{def:saturated}
A closed unital subalgebra $E\subseteq A$ is \emph{saturated} if
$\gauge_E(x)=\gauge_A(x)$ for all $x\in E$, i.e.\ if $E$ retains, for each of
its elements, enough test elements to realise the full ambient budget.
Equivalently, the inclusion $E\hookrightarrow A$ is a PB-morphism.  We call an
isometric PB-morphism a \emph{saturated embedding}.
\end{definition}

\begin{remark}[Correction: monomorphisms]\label{rem:mono-correction}
The earlier notes assert that the monomorphisms of this category are exactly
the gauge-preserving isometric embeddings.  That is not correct as a
categorical statement.  Monomorphism is a cancellation property, and \emph{any}
injective PB-morphism is monic: if $f$ is injective and $fg=fh$ then $g=h$
pointwise.  Nothing about isometry or gauge preservation is needed, and a
contractive injective homomorphism need be neither isometric nor
gauge-preserving.

What is true, and what the earlier notes were reaching for, is that the
saturated embeddings are the right notion of \emph{subobject}: they are the
inclusions along which the ambient structure is faithfully seen, they are
exactly the isometric PB-morphisms by Lemma~\ref{lem:isometric}, and they are
the maps for which the obstruction of \S\ref{sec:categorical} is stated.  The
analogy with $C^*$-algebras is to the fact that an injective
$*$-homomorphism is automatically isometric; the analogue \emph{fails} here,
and that failure is not a defect of the write-up but a real difference between
the theories.
\end{remark}

\begin{remark}[Comparison with $*$-homomorphisms]
On a $C^*$-algebra, $*$-homomorphisms are automatically contractive, and
isometric exactly when injective.  PB-morphisms are budget-contracting by fiat,
and gauge-preserving exactly when isometric.  The structural reason for the
difference is that the $C^*$-norm is computed from data internal to the element
($a$ and $a^*$), whereas the gauge is a supremum over the ambient algebra.  See
Remark~\ref{rem:ambient-referencing}.
\end{remark}

\section{The base functor and the coreflection theorem}
\label{sec:base}

Theorem~\ref{thm:base} attaches a compact Hausdorff space $X_A$ to every object.
This section shows that the assignment is functorial, that it is in fact a
coreflection, and that it exhibits every object as a bundle over its own base.
Everything here is specific to the gauge $\gauge_1$; see
Warning~\ref{warn:vanishing}.

\subsection{Morphisms preserve centres}

\begin{lemma}\label{lem:centre-preserved}
Let $f:A\to B$ be a PB-morphism for the gauge $\gauge_1$.  Then
$f(\Zc(A))\subseteq\Zc(B)$.
\end{lemma}

\begin{proof}
Let $a\in\Zc(A)$.  Then $\gauge_1(a)=0$ by Lemma~\ref{lem:gauge-basic}(e), so
budget-contraction gives $\gauge_1(f(a))\le0$, hence $\ad_{f(a)}=0$, hence
$f(a)\in\Zc(B)$ --- again by Lemma~\ref{lem:gauge-basic}(e), now read in the
other direction.
\end{proof}

\begin{remark}
This is the payoff of the exact identity $\{\gauge_1=0\}=\Zc(A)$.  Note that
the conclusion is genuinely a restriction on arrows and not a triviality: an
ordinary unital contractive homomorphism can certainly move a central element
off the centre, as $\RR^2\hookrightarrow M_2(\RR)$ does.  The morphism law of
Definition~\ref{def:pb-morphism}, introduced to repair the non-functoriality of
the gauge, turns out to buy functoriality of the \emph{base} as well.
\end{remark}

\begin{theorem}[The base is a functor]\label{thm:base-functor}
The assignments $A\mapsto X_A$ and $f\mapsto X_f$ define a functor
\[
  X:\PB^C_{\gauge_1}\longrightarrow\CompHaus\op .
\]
Explicitly, a PB-morphism $f:A\to B$ restricts by
Lemma~\ref{lem:centre-preserved} to a unital homomorphism
$\Zc(f):C(X_A,\RR)\to C(X_B,\RR)$, which by Theorem~\ref{thm:gelfand} is
$X_f^*$ for a unique continuous map $X_f:X_B\to X_A$.  Functoriality is the
functoriality of $\Spec$.
\end{theorem}

\begin{proof}
Lemma~\ref{lem:centre-preserved} gives the restriction, which is unital and
multiplicative; Theorem~\ref{thm:base} identifies source and target with
$C(X_A,\RR)$ and $C(X_B,\RR)$, and Theorem~\ref{thm:gelfand} converts unital
homomorphisms between such algebras into continuous maps of the spectra, in a
way that takes composites to composites and identities to identities.
\end{proof}

\subsection{The commutative part is coreflective}

Write $\CommPB\subseteq\PB^C_{\gauge_1}$ for the full subcategory of
commutative objects.  By Theorem~\ref{thm:gelfand} these are exactly the
algebras $C(X,\RR)$, and the morphisms between them are exactly the unital
homomorphisms --- the budget condition is vacuous, both gauges being zero.
Hence
\[
  \CommPB\;\simeq\;\CompHaus\op .
\]

\begin{theorem}[Coreflection]\label{thm:coreflection}
The centre functor $\Zc:\PB^C_{\gauge_1}\to\CommPB$ is right adjoint to the
inclusion $\iota:\CommPB\hookrightarrow\PB^C_{\gauge_1}$.  That is, for every
commutative object $C$ and every object $A$ there is a bijection, natural in
both variables,
\[
  \Hom_{\PB}\bigl(\iota C,\,A\bigr)\;\cong\;\Hom_{\CommPB}\bigl(C,\,\Zc(A)\bigr).
\]
Equivalently, $\CommPB\simeq\CompHaus\op$ is a coreflective full subcategory of
$\PB^C_{\gauge_1}$, with coreflector $A\mapsto\Zc(A)$.
\end{theorem}

\begin{proof}
First, the inclusion $\Zc(A)\hookrightarrow A$ is a PB-morphism: it is
isometric, and for $z\in\Zc(A)$ both gauges vanish, so the budget condition
holds with equality.  It is therefore a saturated embedding
(Definition~\ref{def:saturated}).

Given a PB-morphism $g:\iota C\to A$ with $C$ commutative, every $c\in C$ has
$\gauge_1^C(c)=0$, so $\gauge_1^A(g(c))=0$ and $g(c)\in\Zc(A)$ by
Lemma~\ref{lem:gauge-basic}(e).  Hence $g$ factors uniquely through
$\Zc(A)\hookrightarrow A$, and the factorisation is a morphism of $\CommPB$
because it is a unital contractive homomorphism between commutative objects.
Conversely, composing a morphism $C\to\Zc(A)$ with the saturated embedding
$\Zc(A)\hookrightarrow A$ gives a PB-morphism $\iota C\to A$.  The two
constructions are mutually inverse, and naturality in $A$ is
Lemma~\ref{lem:centre-preserved} while naturality in $C$ is immediate.
\end{proof}

\begin{corollary}[Points of an object]\label{cor:points}
For every compact Hausdorff $X$ and every object $A$,
\[
  \Hom_{\PB}\bigl(C(X,\RR),\,A\bigr)\;\cong\;C(X_A,X),
\]
the set of continuous maps $X_A\to X$.  In particular $\RR=C(\mathrm{pt},\RR)$
is an initial object, and $\Hom_{\PB}(C(X,\RR),A)$ never sees anything of $A$
beyond its base.
\end{corollary}

\begin{corollary}[Commutative coproducts]\label{cor:comm-coproducts}
The coproduct of $C(X,\RR)$ and $C(Y,\RR)$ in $\PB^C_{\gauge_1}$ exists and
equals $C(X\times Y,\RR)$.
\end{corollary}

\begin{proof}
Let $f:C(X,\RR)\to A$ and $g:C(Y,\RR)\to A$ be PB-morphisms.  By the proof of
Theorem~\ref{thm:coreflection} both land in $\Zc(A)$; in particular their
images commute with one another, so the map $u\otimes v\mapsto f(u)g(v)$ is a
well-defined unital homomorphism on the algebraic tensor product, contractive
for the sup norm, and extends to $C(X\times Y,\RR)=C(X,\RR)\otimes C(Y,\RR)$.
Its image lies in $\Zc(A)$, so its budget vanishes identically and it is a
PB-morphism.  Uniqueness holds because the image of $C(X,\RR)\otimes C(Y,\RR)$
is generated by the two images.
\end{proof}

\begin{warning}[The coproduct is not the free product]
\label{warn:not-free-product}
Corollary~\ref{cor:comm-coproducts} disposes of a suggestion in the earlier
notes, where cocompleteness was to be established by showing that the free
product $A*B$ of PB-algebras is again a PB-algebra.  The free product of
$C(X,\RR)$ and $C(Y,\RR)$ is highly noncommutative, whereas their coproduct in
$\PB$ is the commutative algebra $C(X\times Y,\RR)$.  So the free product is
\emph{not} the coproduct of $\PB$, and the programme ``check that $A*B$ is a
PB-algebra'' does not address cocompleteness.  What Theorem~\ref{thm:coreflection}
shows is that, along the commutative part at least, the coproduct is the one
forced by the geometry: the product of the bases.  The correct general question
is Problem~\ref{prob:coproducts}.
\end{warning}

\subsection{The bundle picture}
\label{sec:bundle}

Theorem~\ref{thm:base} makes $A$ a unital Banach module over $C(X_A,\RR)$, the
action being multiplication by central elements; submultiplicativity gives
$\norm{f\cdot a}\le\norm{f}_\infty\norm{a}$.  This is the standard setting for a
fibrewise decomposition.

\begin{definition}[Fibres]\label{def:fibres}
For $x\in X_A$ let $I_x\subseteq C(X_A,\RR)$ be the ideal of functions
vanishing at $x$, let $J_x=\overline{I_x\cdot A}$ --- a closed two-sided ideal
of $A$, since $I_x\cdot A$ is central-by-$A$ --- and set $A_x=A/J_x$, with
quotient map $\pi_x$.  We call $A_x$ the \emph{fibre of $A$ over $x$}.
\end{definition}

\begin{definition}[Local norm]\label{def:local-norm}
For $a\in A$ and $x\in X_A$ put
\[
  N_x(a)=\inf\bigl\{\norm{f\cdot a}\;:\;f\in C(X_A,\RR),\ 0\le f\le1,\
  f\equiv1\ \text{on a neighbourhood of }x\bigr\}.
\]
\end{definition}

\begin{proposition}\label{prop:usc-field}
For every $a\in A$:
\begin{enumerate}
\item[(a)] $\norm{\pi_x(a)}\le N_x(a)\le\norm{a}$;
\item[(b)] $x\mapsto N_x(a)$ is upper semicontinuous on $X_A$;
\item[(c)] each $\pi_x$ is contractive and budget-contracting, and
  $\pi_x(f\cdot1)=f(x)1$ for $f\in C(X_A,\RR)$.
\end{enumerate}
\end{proposition}

\begin{proof}
(a) If $f\equiv1$ near $x$ then $(1-f)(x)=0$, so $(1-f)\cdot a\in I_x\cdot A\subseteq J_x$
and hence $\pi_x(a)=\pi_x(f\cdot a)$, giving
$\norm{\pi_x(a)}\le\norm{f\cdot a}$; take the infimum.  The upper bound is
$\norm{f\cdot a}\le\norm{f}_\infty\norm{a}\le\norm{a}$.

(b) Given $\varepsilon>0$ choose $f$ admissible for $x$ with
$\norm{f\cdot a}<N_x(a)+\varepsilon$, say $f\equiv1$ on the open set $U\ni x$.
Every $x'\in U$ has $f\equiv1$ on the same neighbourhood $U$ of $x'$, so $f$ is
admissible for $x'$ and $N_{x'}(a)\le\norm{f\cdot a}<N_x(a)+\varepsilon$.

(c) Proposition~\ref{prop:category}, and $f\cdot1-f(x)1\in I_x\cdot A$.
\end{proof}

\begin{remark}[What is missing, precisely]\label{rem:local-norm-gap}
For $C^*$-algebras one has $\norm{\pi_x(a)}=N_x(a)$, and the proof uses an
approximate identity to show that every element of $J_x$ is approximated by
elements $(1-f)\cdot a$.  We do not have that here, and we do not assume it: the
inequality in (a) is all that is claimed.  Consequently
Proposition~\ref{prop:usc-field}(b) gives upper semicontinuity of the
\emph{local} norm $N_x$, and upper semicontinuity of the genuine fibre norm
$x\mapsto\norm{\pi_x(a)}$ follows only once the two agree.
\end{remark}

\begin{question}[Faithfulness of the bundle]\label{q:faithful-bundle}
Is $\norm{a}=\sup_{x\in X_A}\norm{\pi_x(a)}$ for every $a\in A$, and is
$N_x=\norm{\pi_x(\cdot)}$?
\end{question}

For $C^*$-algebras the corresponding statement is a theorem
(Dauns--Hofmann \cite{DaunsHofmann}), and the usual proof needs a
partition-of-unity estimate $\norm{\sum_if_i\cdot a_i}\le\max_i\norm{a_i}$ for
partitions of unity $\{f_i\}$, which holds automatically there but is an extra
hypothesis for a general Banach module --- as is the identification of the
fibre norm with the local norm (Remark~\ref{rem:local-norm-gap}).  A positive answer would give the slogan of this section
its literal meaning:

\begin{quote}
\emph{a PB-algebra is the algebra of sections of an upper semicontinuous field
of pointless PB-algebras over its own classical base.}
\end{quote}

Two ingredients are still missing for the slogan, and both are recorded in
\S\ref{sec:open}: Question~\ref{q:faithful-bundle}, and the question of whether
the fibres are PB-algebras at all, which is the quotient-closure
Problem~\ref{prob:quotients}.  It is worth stressing what is \emph{not}
missing: the base, the central action and the fibring are unconditional, and
they are available because the axioms were built to leave the centre rigid.

\section{Examples and obstructions}
\label{sec:examples}

Throughout this section the gauge is $\gauge_1$ and, unless stated otherwise,
$C=\tfrac14$.

\begin{example}[The commutative objects]
$C(X,\RR)$ for $X$ compact Hausdorff.  Here $\gauge_1\equiv0$, both axioms hold
without slack, $\kap=\lam=0$, and $X_A=X$.
\end{example}

\begin{example}[The quaternions]\label{ex:H}
$\HH$ with its modulus.  The norm is multiplicative, so \eqref{ax:sq} holds
exactly and $\kap(\HH)=0$.  For \eqref{ax:pb2} the extremal element is a purely
imaginary unit $v$: then $v^2=-1$, so
$\bigl\lVert\norm{v^2}1-v^2\bigr\rVert=2$ against $\norm{v^2}=1$, a defect of
$1$, and $\gauge_1(v)=4$ by Proposition~\ref{prop:real-stampfli}.  Hence
$\lam(\HH)=\tfrac14$, attained.  The base is a point: $\HH$ is pointless.
\end{example}

\begin{example}[The $2\times2$ matrices]\label{ex:M2}
$M_2(\RR)$ with the operator norm.  Both constants equal $\tfrac14$ and both are
attained: $\kap$ at $E_{12}$ (Example~\ref{ex:E12}), and $\lam$ at the rotation
$J=\begin{psmallmatrix}0&-1\\1&0\end{psmallmatrix}$, for which $J^2=-1$ gives a
reality defect of $1$ against $\gauge_1(J)=4$.  The base is a point.
\end{example}

\begin{example}[The excluded ones]\label{ex:excluded}
$\CC$ as a real algebra and $\RR[\varepsilon]/(\varepsilon^2)$ are both
commutative, so both have zero budget everywhere, and both have a nonzero
defect: $\CC$ fails \eqref{ax:sr} at $i$, and $\RR[\varepsilon]$ fails
\eqref{ax:sq} at $\varepsilon$.  Neither is a PB-algebra for any constant.  This
pair is the reason the theory is worth stating: it admits noncommutative
``points'' ($M_2(\RR)$, $\HH$) while refusing both complex directions and
infinitesimal ones, and it refuses them for the same reason --- an unfunded
defect.
\end{example}

% Generated by verify/experiments/run_all.py -- do not edit by hand.
\begin{table}[htbp]
\centering
\caption{Critical constants for the gauge $\gauge_1$, found by the
search described in \S\ref{sec:status}.  Each value is a lower bound
for a supremum: \statusnum, not \statusproved.  A dash means the
defect is identically zero; $\infty$ means a nonzero defect with no
budget, so the algebra is not a PB-algebra for any constant.}
\label{tab:constants}
\begin{tabular}{lccc}
\toprule
Algebra & $\kap_{\gauge_1}$ & $\lam_{\gauge_1}$ & $C_{\gauge_1}$ \\
\midrule
$\RR^2$ & --- & --- & 0.0000 \\
$\CC$ & --- & $\infty$ & $\infty$ \\
$\HH$ & --- & 0.2500 & 0.2500 \\
$M_2(\RR)$ & 0.2500 & 0.2500 & 0.2500 \\
\bottomrule
\end{tabular}
\end{table}


\subsection{The two obstructions}

The theory is transverse to the $C^*$ one, and each side has objects the other
does not see.

\subsubsection*{Excess on the PB side: non-central radical}

\begin{proposition}\label{prop:T2}
$T_2(\RR)$, the real upper-triangular $2\times2$ matrices with the operator
norm, has $\rad(T_2(\RR))=\RR E_{12}\ne0$, so its complexification is not a
$C^*$-algebra --- $C^*$-algebras are semisimple.  Its radical is
\emph{non-central}: $[E_{12},\operatorname{diag}(1,-1)]=-2E_{12}\ne0$, so
$\gauge_1(E_{12})=4>0$ and the square defect of $E_{12}$ is funded.
\end{proposition}

Whether $T_2(\RR)$ is an object of $\PB^{1/4}_{\gauge_1}$ --- i.e.\ whether its
funded defect stays within the same constant that $M_2(\RR)$ needs --- is a
quantitative question, settled numerically in the table above and discussed in
\S\ref{sec:status}.  It matters, because $T_2(\RR)$ is the one elementary
finite-dimensional witness for $\PB\not\subseteq\Cstar$, and because it is the
cone that answers Question~\ref{q:legal-cone}.

\subsubsection*{Excess on the $C^*$ side: chirality}

\begin{proposition}\label{prop:chirality}
A $C^*$-algebra $B$ with $B\not\cong\bar B$ admits no real form, hence arises
from no PB-algebra by complexification.  Such $B$ exist, but are necessarily
infinite-dimensional: every finite-dimensional $C^*$-algebra
$\bigoplus_iM_{n_i}(\CC)$ has the real form $\bigoplus_iM_{n_i}(\RR)$.
\end{proposition}

\begin{remark}[Physics]\label{rem:physics}
The structure that actually occurs in physical models is the presence of a
real or quaternionic form: antiunitary time reversal, Dyson's threefold way,
the tenfold way of Altland--Zirnbauer.  That is $B\cong\bar B$ together with a
chosen conjugation --- exactly the PB-side datum.  Genuinely chiral
$B\not\cong\bar B$ arises as a topological invariant of \emph{families} and not
as the observable algebra of a single system: $B(H)$, the CAR and CCR algebras,
AF and UHF algebras and the type III factors all satisfy $B\cong\bar B$.  So the
boundary between the two theories runs parallel to the physical boundary
between systems with and without a real structure.  We record this as
motivation, not as a theorem.
\end{remark}

\subsection{The critical constant of $T_2(\RR)$}
\label{sec:T2-constant}

The earlier notes assert that $\kap_{\gauge_1}(T_2(\RR))=\kap_{\gauge_1}(M_2(\RR))=\tfrac14$,
both attained at the shared element $E_{12}$ with the shared witness
$z=\operatorname{diag}(1,-1)$, and conclude that \emph{no} choice of constant
separates the two algebras.  That conclusion is false.  The extremiser of
$T_2(\RR)$ is not $E_{12}$, and the critical constant is strictly larger than
$\tfrac14$.

The computation is exact, and rests on the fact that commutators in $T_2(\RR)$
span a line.

\begin{lemma}[Dual certificate]\label{lem:dual-certificate}
Let $x=\begin{psmallmatrix}\alpha&\gamma\\0&\beta\end{psmallmatrix}$ and
$z=\begin{psmallmatrix}p&c\\0&q\end{psmallmatrix}$ lie in $T_2(\RR)$.  Then
$[x,z]=L(z)E_{12}$ with $L(z)=c(\alpha-\beta)+\gamma(q-p)$, and for every
$g\in\RR$,
\[
  \gauge_1^{T_2}(x)\;\le\;\bigl\lVert G_g\bigr\rVert_*^2,
  \qquad
  G_g=\begin{pmatrix}-\gamma&\alpha-\beta\\ g&\gamma\end{pmatrix},
\]
where $\norm{\cdot}_*$ is the nuclear norm (the sum of the singular values).
\end{lemma}

\begin{proof}
The formula for $[x,z]$ is a direct computation.  Since $z_{21}=0$, we have
$L(z)=\langle G_g,z\rangle$ in the trace pairing for \emph{every} value of the
free entry $g$.  Duality of the nuclear and operator norms gives
$\abs{\langle G_g,z\rangle}\le\norm{G_g}_*\norm{z}$, so
$\norm{[x,z]}=\abs{L(z)}\le\norm{G_g}_*$ whenever $\norm{z}\le1$.  Take the
supremum over $z$ and then the square.
\end{proof}

\begin{proposition}[$T_2(\RR)$ is not an object at $C=\tfrac14$]
\label{prop:T2-constant}
Let
\[
  \kap^*\;=\;\frac{4y(1-y)}{(1+y)^4}\bigg|_{\,y=\frac{5-\sqrt{17}}{4}}
  \;=\;\frac{128\,(3\sqrt{17}-11)}{15112-3528\sqrt{17}}
  \;=\;0.3098421747\ldots
\]
Then $\kap_{\gauge_1}(T_2(\RR))\ge\kap^*>\tfrac14$.  Consequently
$T_2(\RR)\notin\PB^{1/4}_{\gauge_1}$, and every constant
$C\in[\tfrac14,\kap^*)$ admits $M_2(\RR)$ while excluding $T_2(\RR)$.
\end{proposition}

\begin{proof}
We exhibit a one-parameter family and bound its gauge from above.

For $c\in[0,1]$ put
$x_c=\begin{psmallmatrix}c&1-c^2\\0&-c\end{psmallmatrix}$.  Then
$x_c^\top x_c$ has determinant $\det(x_c)^2=c^4$ and trace $2c^2+(1-c^2)^2$, so
its characteristic polynomial is $\lambda^2-\bigl(2c^2+(1-c^2)^2\bigr)\lambda+c^4$,
which has $\lambda=1$ as a root; the other root is $c^4\le1$.  Hence
$\norm{x_c}=1$.  Also $x_c^2=c^2\cdot1$, so
\[
  \norm{x_c}^2-\norm{x_c^2}=1-c^2 .
\]
Apply Lemma~\ref{lem:dual-certificate} with $\alpha-\beta=2c$ and
$\gamma=1-c^2$.  Write $u=1-c^2$ and $v=2c$ and choose $g=-u^2/v$ (for $c>0$),
which makes
$G_g=\begin{psmallmatrix}-u&v\\-u^2/v&u\end{psmallmatrix}$ have determinant
$-u^2-v\cdot(-u^2/v)=0$.  A rank-one matrix has nuclear norm equal to its
Frobenius norm, so
\[
  \norm{G_g}_*^2=\norm{G_g}_F^2=2u^2+v^2+\frac{u^4}{v^2}
  =\left(\frac{u^2+v^2}{v}\right)^{2}
  =\left(\frac{(1-c^2)^2+4c^2}{2c}\right)^{2}
  =\left(\frac{(1+c^2)^2}{2c}\right)^{2},
\]
using $(1-c^2)^2+4c^2=(1+c^2)^2$.  Therefore
\[
  \gauge_1^{T_2}(x_c)\;\le\;\frac{(1+c^2)^4}{4c^2},
  \qquad\text{hence}\qquad
  \kap_{\gauge_1}(T_2(\RR))\;\ge\;\frac{(1-c^2)\,4c^2}{(1+c^2)^4} .
\]
Writing $y=c^2$ and maximising $f(y)=4y(1-y)/(1+y)^4$ on $(0,1)$: the
stationarity condition $(1-2y)(1+y)-4y(1-y)=0$ reduces to $2y^2-5y+1=0$, whose
root in $(0,1)$ is $y=(5-\sqrt{17})/4=0.2192235\ldots$, giving
$f(y)=\kap^*$.  Finally $\kap^*>\tfrac14$ by direct evaluation.
\end{proof}

\begin{remark}[What went wrong in the earlier argument]
\label{rem:T2-correction}
The earlier notes located the extremiser of $T_2(\RR)$ at $E_{12}$, by analogy
with $M_2(\RR)$ and on the grounds that ``scalar diagonal minimises $\gauge$ for
fixed nilpotent content''.  The optimum is instead at the traceless element
$x_c$ with $c=\sqrt{(5-\sqrt{17})/4}=0.46821\ldots$, which is neither nilpotent
nor a perturbation of one: it has $x_c^2=c^2\cdot1$, a square defect of
$1-c^2=0.7808\ldots$, and a budget of only $2.5199\ldots$ rather than the $4$
available at $E_{12}$.  The trade is that moving away from $E_{12}$ costs some
defect but costs \emph{more} budget, because the two diagonal entries pull the
optimal test element in competing directions.  This is exactly the sort of claim
that a numerical sweep settles and an analogy does not.
\end{remark}

\begin{remark}[Numerical corroboration]
An independent search over all of $T_2(\RR)$, with the gauge computed by the
convex dual of Lemma~\ref{lem:dual-certificate}, finds
$\kap_{\gauge_1}(T_2(\RR))=0.30984217\ldots$ and
$\lam_{\gauge_1}(T_2(\RR))=0.11294508\ldots$, so that
$C_{\gauge_1}(T_2(\RR))=\kap^*$ and the bound of
Proposition~\ref{prop:T2-constant} is attained.  That the supremum over all of
$T_2(\RR)$ is attained on the traceless slice is \statusnum; the inequality
$\kap\ge\kap^*$, which is what separates $T_2(\RR)$ from $M_2(\RR)$, is
\statusproved.  See \texttt{verify/experiments/t2\_exact.py}.
\end{remark}

\section{Categorical status}
\label{sec:categorical}

We work in $\PB^C_{\gauge_1}$ for a fixed $C$, and we assume
Convention~\ref{conv:degenerate}.

\subsection{Initial and terminal objects}

\begin{proposition}\label{prop:terminal}
$\RR$ is initial and the degenerate algebra $0$ is terminal.  Without
Convention~\ref{conv:degenerate}, and for $C\ge\tfrac14$, there is no terminal
object and hence no limits --- for reasons unrelated to the subject.
\end{proposition}

\begin{proof}
For initiality, the unit $\RR\to A$, $\lambda\mapsto\lambda1$, is the unique
unital homomorphism; it is isometric and both gauges vanish on its image.  For
terminality, $A\to0$ is unique, contractive and budget-contracting.

Without the degenerate object, suppose $T$ were terminal.  Since
$C\ge\tfrac14$, $M_2(\RR)$ is an object (Example~\ref{ex:M2}); let
$\varphi:M_2(\RR)\to T$ be the unique morphism.  As $M_2(\RR)$ is simple and
$\varphi$ is unital, $\varphi$ is injective.  For $v\in O(2)$ the automorphism
$\Ad_v$ of $M_2(\RR)$ is isometric, hence gauge-preserving by
Lemma~\ref{lem:isometric}, hence a PB-morphism; so $\varphi\circ\Ad_v$ is again
a morphism $M_2(\RR)\to T$, and uniqueness gives
$\varphi\circ\Ad_v=\varphi$.  Injectivity of $\varphi$ then forces
$\Ad_v=\id$ for every $v\in O(2)$, which is false.  Hence there is no terminal
object, and since the terminal object is the empty product, no limits.  (The
same argument runs with $\HH$ and its inner automorphisms in place of
$M_2(\RR)$; some noncommutative object is needed, which is why the statement is
restricted to $C\ge\tfrac14$.)
\end{proof}

\subsection{Products}

\begin{proposition}\label{prop:products}
$\PB^C_{\gauge_1}$ has all products.  For $A=\prod_iA_i$ with the sup norm and
$a=(a_i)$,
\[
  \gauge_1(a)=\sup_i\gauge_1(a_i),
\]
and the projections are PB-morphisms with the expected universal property.
\end{proposition}

\begin{proof}
For the gauge formula, $\norm{[a,z]}=\sup_i\norm{[a_i,z_i]}$ and $\norm{z}\le1$
forces $\norm{z_i}\le1$, giving $\le$; conversely a test element concentrated in
one factor gives $\ge$.

For \eqref{ax:pb1} let $M=\norm{a}=\sup_i\norm{a_i}$ and fix $\varepsilon>0$ and
an index $i$ with $\norm{a_i}\ge M-\varepsilon$.  Then
\[
  \norm{a}^2-\norm{a^2}
  \;\le\;M^2-\norm{a_i^2}
  \;=\;\bigl(M^2-\norm{a_i}^2\bigr)+\bigl(\norm{a_i}^2-\norm{a_i^2}\bigr)
  \;\le\;2M\varepsilon+C\gauge_1(a_i)
  \;\le\;2M\varepsilon+C\gauge_1(a),
\]
and $\varepsilon\downarrow0$ gives the claim.  For \eqref{ax:pb2} write
$S=\norm{a^2}=\sup_j\norm{a_j^2}$; then for each $i$,
\[
  \bigl\lVert S1-a_i^2\bigr\rVert
  \;\le\;\bigl\lVert\,\norm{a_i^2}1-a_i^2\,\bigr\rVert+\bigl(S-\norm{a_i^2}\bigr)
  \;\le\;\norm{a_i^2}+C\gauge_1(a_i)+S-\norm{a_i^2}
  \;\le\;S+C\gauge_1(a),
\]
and taking the supremum over $i$ on the left gives
$\bigl\lVert S1-a^2\bigr\rVert\le S+C\gauge_1(a)$.  Projections contract both norm and gauge; for
a cone $(f_i:T\to A_i)$ the map $f=(f_i)$ satisfies
$\gauge_1(f(x))=\sup_i\gauge_1(f_i(x))\le\gauge_1(x)$ and is contractive, and it
is the unique such factorisation.
\end{proof}

\begin{remark}[Uniformity of the constant is essential]
The proof uses one constant for all factors.  In $\PB_{\gauge_1}=\bigcup_C\PB^C$
products need not exist: a family with $C_{\gauge_1}(A_i)\to\infty$ has a product
whose critical constant is the supremum, hence infinite.  This is a reason to
regard the filtered pieces $\PB^C$, and not their union, as the categories of
interest.
\end{remark}

\subsection{Subalgebras, and why equalizers are hard}

\begin{proposition}[The class is not closed under subalgebras]
\label{prop:not-closed-subalgebras}
A closed unital subalgebra of a PB-algebra need not be a PB-algebra.  Indeed
$\RR1\oplus\RR E_{12}\cong\RR[\varepsilon]/(\varepsilon^2)$ is a closed unital
subalgebra of $M_2(\RR)$, and by Example~\ref{ex:excluded} it is not a
PB-algebra for any constant.
\end{proposition}

\begin{remark}[Ambient referencing]\label{rem:ambient-referencing}
Proposition~\ref{prop:not-closed-subalgebras} isolates the structural
difficulty.  The gauge is defined by a supremum over test elements $z$ drawn
from the ambient algebra, so passing to a subalgebra can destroy the budget
while preserving the defect that the budget was paying for.  A $C^*$-algebra
has no analogous failure: the datum defining the norm, $r(a^*a)^{1/2}$, is built
from $a$ and $a^*$, both of which lie in any $*$-subalgebra containing $a$.
Closure under $*$-subalgebras is free there, and that is exactly why $C^*$
has equalizers --- the solution set $\{f=g\}$ is a norm-closed $*$-subalgebra,
hence an object, and its inclusion is isometric, hence legal.  Everything
peculiar about limits in $\PB$ traces back to this one asymmetry.
\end{remark}

\begin{proposition}[The set-theoretic equalizer is not legal]
\label{prop:delta-not-saturated}
Let $A=M_2(\RR)$, $u=\operatorname{diag}(1,-1)$, and consider the parallel pair
$\id,\Ad_u:A\rightrightarrows A$.  The set-theoretic equalizer is the diagonal
$\Delta=\{\operatorname{diag}(\alpha,\beta)\}\cong\RR^2$, which \emph{is} a
PB-algebra, but its inclusion $\Delta\hookrightarrow M_2(\RR)$ is \emph{not} a
PB-morphism: $\gauge_1^\Delta\equiv0$ while
$\gauge_1^{M_2}(\operatorname{diag}(\alpha,\beta))=(\alpha-\beta)^2$, so
$\Delta$ is not saturated.  The largest saturated subalgebra of $\Delta$ is
$\RR1$.
\end{proposition}

\begin{proof}
$\Delta$ is commutative, so its own gauge vanishes.  In $M_2(\RR)$,
$[\operatorname{diag}(\alpha,\beta),z]$ has entries
$(\alpha-\beta)z_{12}$ and $-(\alpha-\beta)z_{21}$ off the diagonal and zero on
it, so its norm is $\abs{\alpha-\beta}\max(\abs{z_{12}},\abs{z_{21}})\le\abs{\alpha-\beta}$
with equality for $z=E_{12}$; this also agrees with
Proposition~\ref{prop:real-stampfli}, since
$\dist(\operatorname{diag}(\alpha,\beta),\RR1)=\tfrac12\abs{\alpha-\beta}$.  The
subalgebras of $\Delta\cong\RR^2$ are $\RR1$ and $\Delta$, and $\RR1$ is
saturated because both gauges vanish on it.
\end{proof}

\begin{remark}[A retraction, upheld]
The earlier notes first concluded from Proposition~\ref{prop:delta-not-saturated}
that $\PB$ is incomplete, then correctly retracted the argument: a categorical
equalizer need not be the set-theoretic one, and the natural candidate is
instead the largest saturated subalgebra $\RR1$, whose universality has to be
tested against \emph{legal} cones only.  That retraction stands.  What follows
settles the question it left behind.
\end{remark}

\subsection{Legal cones: the question, and its answer}

\begin{question}[\cite{notes}, Question~19]\label{q:legal-cone}
Does there exist an object $T$ and a PB-morphism $t:T\to M_2(\RR)$ with
$t(T)\subseteq\Delta$ and $t(T)\ne\RR1$?
\end{question}

The question has a clean reformulation, which also explains why it is not
obviously answerable either way.

\begin{lemma}[Cones are admissible character pairs]\label{lem:cones-characters}
PB-morphisms $t:T\to M_2(\RR)$ with $t(T)\subseteq\Delta$ correspond bijectively
to ordered pairs $(\chi_1,\chi_2)$ of real characters of $T$ satisfying
\begin{equation}\label{eq:admissible}
  \bigl(\chi_1(x)-\chi_2(x)\bigr)^2\;\le\;\gauge_1^T(x)\qquad\text{for all }x\in T,
\end{equation}
via $t(x)=\operatorname{diag}(\chi_1(x),\chi_2(x))$.  Every such $t$ equalizes
$\id$ and $\Ad_u$.  Call a pair satisfying \eqref{eq:admissible}
\emph{admissible}.
\end{lemma}

\begin{proof}
A unital homomorphism into $\Delta\cong\RR^2$ is a pair of unital
homomorphisms $T\to\RR$, i.e.\ a pair of real characters.  Contractivity is
automatic: $\norm{\operatorname{diag}(\chi_1(x),\chi_2(x))}=\max_i\abs{\chi_i(x)}\le r(x)\le\norm{x}$.
The budget condition is $\gauge_1^{M_2}(t(x))\le\gauge_1^T(x)$, and
$\gauge_1^{M_2}(\operatorname{diag}(\alpha,\beta))=(\alpha-\beta)^2$ by
Proposition~\ref{prop:delta-not-saturated}.  Finally $\Ad_u$ fixes $\Delta$
pointwise, so any such $t$ equalizes the pair.
\end{proof}

\begin{lemma}[Admissible pairs agree on the centre]\label{lem:agree-on-centre}
If $(\chi_1,\chi_2)$ is admissible then $\chi_1=\chi_2$ on $\Zc(T)$.
Equivalently, the two characters lie over the same point of the classical base
$X_T$.
\end{lemma}

\begin{proof}
For $x\in\Zc(T)$ we have $\gauge_1^T(x)=0$, so \eqref{eq:admissible} forces
$\chi_1(x)=\chi_2(x)$.
\end{proof}

So a legal cone with non-scalar image requires an object with two distinct
characters in a single fibre of its base --- a pointless-in-the-relevant-fibre
object that is noncommutative enough to fund the separation of its own
characters.  The commutative objects cannot do it, by
Lemma~\ref{lem:agree-on-centre}; $M_2(\RR)$ cannot do it, having no characters
at all.  But there is an object that can.

\begin{theorem}[Answer to Question~\ref{q:legal-cone}, for $C\ge\kap^*$]
\label{thm:t2-cone}
Let $C\ge\kap^*=0.3098\ldots$, so that $T=T_2(\RR)$ is an object of
$\PB^C_{\gauge_1}$ (Proposition~\ref{prop:T2-constant}), and let
\[
  t:T_2(\RR)\longrightarrow M_2(\RR),
  \qquad
  t\begin{pmatrix}\alpha&\gamma\\0&\beta\end{pmatrix}
  =\begin{pmatrix}\alpha&0\\0&\beta\end{pmatrix}
\]
be the map killing the radical.  Then $t$ is a PB-morphism, it equalizes $\id$
and $\Ad_u$, and its image is all of $\Delta$.  Consequently $\RR1$ is
\emph{not} the equalizer of $\id$ and $\Ad_u$ in $\PB^C_{\gauge_1}$.
\end{theorem}

\begin{proof}
$t$ is the quotient by $\rad(T_2(\RR))=\RR E_{12}$ followed by the identification
$T_2(\RR)/\RR E_{12}\cong\RR^2\cong\Delta$, so it is a unital homomorphism with
image $\Delta$.  Its two components are the characters
$\chi_1(x)=\alpha$ and $\chi_2(x)=\beta$; by Lemma~\ref{lem:cones-characters} it
remains to verify admissibility \eqref{eq:admissible}.

Let $x=\begin{psmallmatrix}\alpha&\gamma\\0&\beta\end{psmallmatrix}$ and
$z=\begin{psmallmatrix}p&c\\0&q\end{psmallmatrix}$.  A direct computation gives
\[
  [x,z]=\bigl(c(\alpha-\beta)+\gamma(q-p)\bigr)E_{12},
\]
so taking $z=E_{12}$, which lies in $T_2(\RR)$ and has norm $1$, yields
$\norm{[x,E_{12}]}=\abs{\alpha-\beta}$ and therefore
\[
  \gauge_1^{T_2}(x)\;\ge\;(\alpha-\beta)^2=\bigl(\chi_1(x)-\chi_2(x)\bigr)^2 .
\]
That is exactly \eqref{eq:admissible}.  Finally, if $\RR1$ with its inclusion
were the equalizer, the cone $t$ would factor through it, forcing
$t(T_2(\RR))\subseteq\RR1$, contradicting $t(T_2(\RR))=\Delta$.
\end{proof}

\begin{remark}[The mechanism]
$T_2(\RR)$ has two characters, the two diagonal entries, separated by a
non-central radical direction; that direction is precisely what funds the budget
realising their separation.  In the language of \S\ref{sec:visibility} it is an
object whose noncommutativity is Jacobson-invisible but commutator-visible ---
and it is the commutator that the morphism law measures.
\end{remark}

\begin{warning}[The answer depends on the constant]\label{warn:cone-constant}
Theorem~\ref{thm:t2-cone} requires $T_2(\RR)$ to be an \emph{object}, and by
Proposition~\ref{prop:T2-constant} it is one only for $C\ge\kap^*>\tfrac14$.  So
the status of Question~\ref{q:legal-cone} is not uniform in $C$:
\begin{itemize}
\item for $C\ge\kap^*$ the answer is \emph{yes} and $\RR1$ is not the equalizer;
\item for $C\in[\tfrac14,\kap^*)$ --- in particular at the critical constant
  $C=\tfrac14$, where $M_2(\RR)$ lives and $T_2(\RR)$ does not --- the question
  remains \emph{open}: the cone constructed above is unavailable, and we know of
  no other.
\end{itemize}
This is worth stating plainly because it is easy to miss.  A cone is an object
of the category, so an answer to a limit question can change as the constant
crosses a threshold, and $\kap^*$ is such a threshold.  Note also what
Lemma~\ref{lem:agree-on-centre} demands of any cone at $C=\tfrac14$: an object
with two distinct characters lying over one point of its base, with their
separation funded by its own commutators.  If
Conjecture~\ref{conj:semisimple} holds, such an object is semisimple, which
makes the search considerably harder --- the obvious source of two characters
in one fibre is a non-central radical, and that is exactly what
semisimplicity forbids.
\end{warning}

\begin{question}[Does the equalizer exist?]\label{q:equalizer-exists}
Theorem~\ref{thm:t2-cone} eliminates the natural candidate but does not decide
existence.  By Lemma~\ref{lem:cones-characters} the equalizer of $\id$ and
$\Ad_u$ exists if and only if the functor
\[
  F:\PB^C_{\gauge_1}\longrightarrow\cat{Set},
  \qquad
  F(T)=\bigl\{\text{admissible pairs of characters of }T\bigr\}
\]
is representable, in which case the representing object, equipped with its
universal admissible pair, is the equalizer.
\end{question}

This is a sharper question than the original, and it is the form in which we
recommend attacking completeness.  Two remarks on what a solution must handle.
$F(M_2(\RR))=\varnothing$, so a representing object $E$ admits no morphism from
$M_2(\RR)$; and $F$ takes a commutative $C(X,\RR)$ to the diagonal pairs only,
so $\Hom(C(X,\RR),E)$ must be the set of continuous maps $X_E\to X$ by
Corollary~\ref{cor:points} --- a strong constraint on the base of $E$.

\subsection{Colimits}

\begin{proposition}\label{prop:comm-colimits}
By Theorem~\ref{thm:coreflection} the inclusion
$\CommPB\hookrightarrow\PB^C_{\gauge_1}$ is a left adjoint, hence preserves all
colimits that exist in $\CommPB\simeq\CompHaus\op$.  In particular coproducts of
commutative objects exist and are given by products of bases
(Corollary~\ref{cor:comm-coproducts}).
\end{proposition}

\begin{remark}[Status of cocompleteness]
Warning~\ref{warn:not-free-product} disposes of the programme of the earlier
notes, which proposed to establish cocompleteness by checking that free products
of PB-algebras are PB-algebras: the free product is not the coproduct here.
Coequalizers reduce to the question of whether $\PB$ is closed under quotients
by closed two-sided ideals --- quotient maps are always legal by
Proposition~\ref{prop:category}, so the only issue is whether the quotient
satisfies the axioms --- which is Problem~\ref{prob:quotients}, and which is
also exactly what the bundle picture of \S\ref{sec:bundle} needs.  We regard
``cocomplete but not complete'' as the expected outcome, and note for the record
that the earlier notes' retraction of the reflector argument was correct: a
reflective subcategory of a complete category inherits limits, so a reflector
$\BanAlg\to\PB$ would force completeness.
\end{remark}

\section{Comparison with $C^*$-algebras}
\label{sec:cstar}

\subsection{The two embeddings of classical geometry}

Both theories extend the same commutative picture, along different embeddings:
\[
  F_{\Cstar}=C(-,\CC),\qquad F_{\PB}=C(-,\RR),
  \qquad\text{related by }\ \Lambda=(-)\otimes_\RR\CC,\quad \Lambda\circ F_{\PB}\cong F_{\Cstar}.
\]
Any comparison must be taken relative to $\Lambda$: as a \emph{real} algebra,
$C(X,\CC)$ is not a PB-algebra, since it contains the central element $i$ and
fails \eqref{ax:pb2} there (Example~\ref{ex:excluded}).

\begin{definition}\label{def:intersection}
Let $\mathcal{I}$ be the class of PB-algebras $A$ whose complexification
$A_\CC$ carries a compatible $C^*$-structure --- equivalently, the real forms of
$C^*$-algebras that satisfy \eqref{ax:pb1} and \eqref{ax:pb2}.
\end{definition}

$\mathcal{I}$ is populated: $M_2(\RR)$ and $\HH$ both complexify to $M_2(\CC)$,
and both are PB-algebras with constant $\tfrac14$ (Examples~\ref{ex:H},
\ref{ex:M2}).  Note that these are two different real forms of one
$C^*$-algebra, so $\Lambda$ is far from injective on objects.  The relationship
between the theories is a span
\[
  \PB\;\hookleftarrow\;\mathcal{I}\;\longrightarrow\;\Cstar,
\]
and not an inclusion in either direction --- with a caveat on the first half
that the earlier notes do not record.  Proposition~\ref{prop:chirality} gives
$C^*$-algebras outside the image (no real form at all), at every constant.  In
the other direction, Proposition~\ref{prop:T2} gives a PB-object outside
$\mathcal{I}$ --- $T_2(\RR)$ has non-central radical, impossible in a semisimple
$C^*$-algebra --- but only for $C\ge\kap^*$, since below that $T_2(\RR)$ is not
an object (Proposition~\ref{prop:T2-constant}).  Whether
$\PB^{1/4}_{\gauge_1}\subseteq\mathcal{I}$, i.e.\ whether the span degenerates to
an inclusion at the critical constant, is open and is implied by
Conjecture~\ref{conj:semisimple} only in part: semisimplicity is necessary for
membership in $\mathcal{I}$, not sufficient.

\subsection{Forgetting the involution}

\begin{proposition}\label{prop:U-faithful}
Let $U$ be the forgetful assignment taking a real form of a $C^*$-algebra, with
its $*$-homomorphisms, to the underlying PB-algebra with its PB-morphisms.
Then $U$ is faithful and bijective on objects but \emph{not full}: there are
PB-morphisms between objects of $\mathcal{I}$ that are not induced by
$*$-homomorphisms.
\end{proposition}

\begin{proof}[Proof sketch]
Faithfulness and bijectivity on objects are immediate.  For fullness, take
$A=M_2(\RR)$ and $g\in GL_2(\RR)$ with $\Ad_g$ isometric and gauge-preserving
but $g$ not a scalar multiple of an orthogonal matrix; then $\Ad_g$ is a
PB-automorphism of $A$ whose complexification is not a $*$-automorphism of
$M_2(\CC)$, since $*$-automorphisms of $M_2(\CC)$ are exactly the $\Ad_v$ with
$v$ unitary.
\end{proof}

\begin{remark}
So on the intersection the PB-theory has \emph{strictly more} morphisms off the
centre, and $\mathcal{I}$ carries two different subobject lattices:
$*$-subalgebras on one side, saturated subalgebras
(Definition~\ref{def:saturated}) on the other.  $U$ is an equivalence exactly on
the commutative floor, where the gauge vanishes identically and every unital
homomorphism is legal.  $\mathcal{I}$ is a span-vertex, not a subcategory of
either theory in a way that respects both structures --- and in particular there
is no reason for it to be closed under PB-quotients.
\end{remark}

\subsection{A uniform bound on the $C^*$ side}

The examples suggest something stronger than membership case by case.

\begin{conjecture}[Uniform bound]\label{conj:cstar-uniform}
Every real $C^*$-algebra $A$ --- a real form of a $C^*$-algebra with its
inherited norm --- satisfies \eqref{ax:pb1} with gauge $\gauge_1$ and constant
$C=\tfrac14$.  Equivalently, by Proposition~\ref{prop:real-stampfli},
\[
  \norm{a}^2-\norm{a^2}\;\le\;\dist(a,\RR1)^2\qquad\text{for all }a .
\]
\end{conjecture}

\begin{remark}
The evidence is as follows.  Equality holds at every square-zero element: if
$a^2=0$ and $\norm{a}=1$ then the left side is $1$, and writing $a$ in its
block form shows $\norm{a-t1}\ge\norm{a}$ for every $t$, so the right side is
$1$ as well.  The inequality also holds at every normal element, where the left
side vanishes.  Numerically it survives a search over $M_n(\RR)$ for
$n\le4$ (\S\ref{sec:status}).  A proof would be valuable well beyond this
subject: it says that failure of the square property is controlled by distance
to the scalars, uniformly, with the sharp constant attained on nilpotents.
Note that \eqref{ax:pb2} is a genuinely separate matter and is \emph{not} implied:
$A=\CC$ is a real $C^*$-algebra failing \eqref{ax:pb2}, and so is $M_2(\CC)$
viewed as a real algebra, whose centre $\CC$ is not of the form $C(X,\RR)$.  The
gate is always the centre.
\end{remark}

\section{Two notions of visibility}
\label{sec:visibility}

The examples separate two senses in which a ``nilpotent direction'' can be
visible, and keeping them apart explains exactly which algebras the two theories
admit.

\begin{definition}\label{def:visibility}
Let $a\in A$.
\begin{itemize}
\item $a$ is \emph{Jacobson-invisible} if $a\in\rad(A)$, i.e.\ $a$ is killed by
  every irreducible representation.  This is the scheme-theoretic sense of an
  infinitesimal thickening, as in $\RR[\varepsilon]/(\varepsilon^2)$.
\item $a$ is \emph{commutator-invisible} if $\gauge_1(a)=0$, i.e.\
  $a\in\Zc(A)$.
\end{itemize}
\end{definition}

The two agree on central elements (both invisible) and on matrix nilpotents in a
simple block (both visible), but they diverge in general.

\begin{proposition}[$T_2(\RR)$ separates the notions]\label{prop:visibility-T2}
In $T_2(\RR)$ the element $E_{12}$ lies in the radical, hence is
Jacobson-invisible --- it is killed by both diagonal characters --- yet
$\gauge_1(E_{12})=4>0$, so it is commutator-visible.  Consequently
$\gauge_1$ tracks non-centrality, which coincides with representational
visibility only when $\rad(A)\subseteq\Zc(A)$.
\end{proposition}

\begin{remark}[A claim that needs the constant]\label{rem:excess}
The earlier notes summarise the situation as: the excess of $\PB$ over the
intersection $\mathcal{I}$ is exactly the non-central radical, because the
$C^*$-instrument $a\mapsto a^*a$ forbids radical outright while $\gauge_1$
forbids only \emph{central} radical.  The first half of that is right --- central
radical is unfunded, hence always excluded.  The second half is not a theorem,
and at $C=\tfrac14$ it is contradicted by the one case we can compute:
$T_2(\RR)$ has non-central radical and is \emph{not} an object at $\tfrac14$
(Proposition~\ref{prop:T2-constant}).  What is true is that $\gauge_1$ forbids
central radical at every constant, and forbids the non-central radical of
$T_2(\RR)$ below $\kap^*$.  Whether it forbids \emph{all} radical at
$C=\tfrac14$ is Conjecture~\ref{conj:semisimple}.
\end{remark}

\begin{proposition}[$\gauge_\infty$ discriminates but is fragile]
\label{prop:dinf-fragile}
The spectral gauge excludes $T_2(\RR)$: every commutator there is a multiple of
the nilpotent $E_{12}$, so $\gauge_\infty\equiv0$ on $T_2(\RR)$ while the square
defect of $E_{12}$ is $1$, which is therefore unfunded.  It admits $M_2(\RR)$,
where $\gauge_\infty(E_{12})=1$ via the non-nilpotent commutator
$[E_{12},E_{21}]=E_{11}-E_{22}$.  This is a genuinely spectral,
semisimplicity-sensitive cut, matching the $C^*$ exclusion of the radical.  But
$\gauge_\infty$ vanishes on whole radical-commutator directions, and quotients
can manufacture such directions, so a $\gauge_\infty$-class has no reason to be
closed under quotients.
\end{proposition}

\begin{remark}[The trade-off, tabulated]
\begin{center}
\begin{tabular}{lll}
\toprule
Gauge & Binding budget on $T_2(\RR)$ & Behaviour \\
\midrule
$\gauge_1$ & $\gauge_1(E_{12})=4$, defect funded (but $\kap=\kap^*>\tfrac14$) & closure-friendly; base is functorial \\
large finite $q$ & $\gauge_q$ small but nonzero & plausibly closed; $T_2$ approximately out \\
$\gauge_\infty$ & $\gauge_\infty\equiv0$, defect unfunded & $T_2$ excluded; closure fails \\
\bottomrule
\end{tabular}
\end{center}
The first row's parenthetical is where the earlier notes and the present
treatment part company: at $C=\tfrac14$ the norm gauge \emph{does} exclude
$T_2(\RR)$ (Theorem~\ref{thm:constant-separates}), so the first row buys
discrimination as well as closure, and the trade-off the table describes may be
an artefact of looking only at constants above $\kap^*$.  The middle row is the
hope that a single large finite $q$ gets the discrimination of $\gauge_\infty$
without its zeros.  It is
Problem~\ref{prob:finite-q}, and we stress that
Lemma~\ref{lem:gauge-basic}(e) makes it \emph{not} merely a matter of tuning:
already at $q=2$ the vanishing locus is strictly larger than the centre, so the
functoriality of \S\ref{sec:base} is lost at every finite $q\ge2$, not only in
the limit.  Any positive answer to Problem~\ref{prob:finite-q} therefore buys
discrimination at the price of the geometry.
\end{remark}

\section{Status ledger}
\label{sec:status}

The table records the epistemic status of every substantive claim above.
\statusproved\ means a proof is given here or a precise citation is supplied;
\statusnum\ means a finite-dimensional inequality checked by the code in
\texttt{verify/}, not proved; \statusconj\ and \statusopen\ mean what they say.

\begin{center}
\begin{tabular}{p{0.66\textwidth}l}
\toprule
Statement & Status \\
\midrule
\eqref{ax:sq}$+$\eqref{ax:sr} $\Rightarrow$ $A\cong C(X,\RR)$, commutativity forced
  (Thms.~\ref{thm:gelfand}, \ref{thm:commutativity-forced})
  & \statusproved$^{\dagger}$ \\
$\Zc(A)\cong C(X_A,\RR)$ for every PB-algebra (Thm.~\ref{thm:base})
  & \statusproved \\
Characters are approximately real, $(\Ima\chi(a))^2\le C\gauge_1(a)$
  (Thm.~\ref{thm:approx-real}) & \statusproved \\
$r(a)\ge(1-4C)\norm{a}$; subcritical objects are semisimple and nilpotent-free
  (Thm.~\ref{thm:normaloid}, Cor.~\ref{cor:subcritical-rigid}) & \statusproved \\
Finite-dimensional zero-budget survivors are products of $\RR$ and $\HH$
  (Prop.~\ref{prop:strict-survivors}) & \statusproved \\
$\gauge_1$ vanishes exactly on the centre; $\gauge_q$, $q\ge2$, does not
  (Lem.~\ref{lem:gauge-basic}(e), Warning~\ref{warn:vanishing}) & \statusproved \\
Isometric PB-morphisms preserve the gauge (Lem.~\ref{lem:isometric})
  & \statusproved \\
PB-morphisms preserve centres; $X$ is a functor
  (Lem.~\ref{lem:centre-preserved}, Thm.~\ref{thm:base-functor}) & \statusproved \\
$\CommPB\simeq\CompHaus\op$ is coreflective in $\PB$ (Thm.~\ref{thm:coreflection})
  & \statusproved \\
Coproducts of commutative objects are $C(X\times Y,\RR)$; the free product is
  \emph{not} the coproduct (Cor.~\ref{cor:comm-coproducts}) & \statusproved \\
$\PB^C_{\gauge_1}$ has all products; the union over $C$ does not
  (Prop.~\ref{prop:products}) & \statusproved \\
No terminal object without the degenerate algebra (Prop.~\ref{prop:terminal})
  & \statusproved \\
$\PB$ is not closed under closed subalgebras
  (Prop.~\ref{prop:not-closed-subalgebras}) & \statusproved \\
$\Delta\hookrightarrow M_2(\RR)$ is not legal; $\RR1$ is the largest saturated
  subalgebra of $\Delta$ (Prop.~\ref{prop:delta-not-saturated}) & \statusproved \\
A legal cone onto all of $\Delta$ exists for $C\ge\kap^*$, so $\RR1$ is not the
  equalizer there (Thm.~\ref{thm:t2-cone}) & \statusproved \\
The same question at $C\in[\tfrac14,\kap^*)$ (Warning~\ref{warn:cone-constant})
  & \statusopen \\
$\kap_{\gauge_1}(T_2(\RR))\ge\kap^*>\tfrac14$, so a constant separates
  $T_2(\RR)$ from $M_2(\RR)$ (Prop.~\ref{prop:T2-constant},
  Thm.~\ref{thm:constant-separates}) & \statusproved \\
Monomorphisms are the gauge-preserving embeddings & \emph{withdrawn}
  (Rem.~\ref{rem:mono-correction}) \\
Every $\gauge_q$ vanishes exactly on the centre & \emph{withdrawn}
  (Warning~\ref{warn:vanishing}) \\
No $\gauge_1$-constant separates $T_2(\RR)$ from $M_2(\RR)$; both have
  $\kap=\tfrac14$ at $E_{12}$ & \emph{refuted}
  (Rem.~\ref{rem:separation-correction}) \\
The excess of $\PB$ over $\mathcal{I}$ is exactly non-central radical
  & \emph{qualified} (Rem.~\ref{rem:excess}) \\
$T_2(\RR)$ has non-central radical, is outside $\mathcal{I}$
  (Props.~\ref{prop:T2}, \ref{prop:visibility-T2}) & \statusproved \\
$\gauge_\infty$ excludes $T_2(\RR)$ and admits $M_2(\RR)$
  (Prop.~\ref{prop:dinf-fragile}) & \statusproved \\
$U:\mathcal{I}_{\Cstar}\to\mathcal{I}_{\PB}$ faithful, bijective on objects, not
  full (Prop.~\ref{prop:U-faithful}) & \statusproved \\
Chirally asymmetric $C^*$-algebras have no real form
  (Prop.~\ref{prop:chirality}) & \statusproved \\
\midrule
$\gauge_1(a)=4\dist(a,\RR1)^2$ on $M_n(\RR)$, $n\le4$
  (Prop.~\ref{prop:real-stampfli}) & \statusnum \\
The identity fails on $T_n(\RR)$ (Prop.~\ref{prop:stampfli-fails-T}) & \statusnum \\
$\kap=\lam=\tfrac14$ for $M_2(\RR)$; $\kap=0$, $\lam=\tfrac14$ for $\HH$
  & \statusnum \\
$\kap_{\gauge_1}(T_2(\RR))=\kap^*$ exactly (the traceless slice is extremal)
  & \statusnum \\
$\kap_{\gauge_1}(M_n(\RR))=\tfrac14$ for $n\le4$
  (Prob.~\ref{prob:matrix-stability}) & \statusnum \\
\midrule
Uniform $C^*$ bound $\norm{a}^2-\norm{a^2}\le\dist(a,\RR1)^2$
  (Conj.~\ref{conj:cstar-uniform}) & \statusconj \\
Semisimplicity at $C=\tfrac14$ (Conj.~\ref{conj:semisimple}) & \statusconj \\
Quaternionic Stone--Weierstrass (Prob.~\ref{prob:quaternionic-sw}) & \statusconj \\
Completeness / representability of $F$ (Prob.~\ref{prob:completeness}) & \statusopen \\
Closure under quotients (Prob.~\ref{prob:quotients}) & \statusopen \\
Coproducts, hence cocompleteness (Prob.~\ref{prob:coproducts}) & \statusopen \\
Faithfulness of the bundle (Prob.~\ref{prob:faithful}) & \statusopen \\
Existence of an intrinsic budget (Prob.~\ref{prob:intrinsic}) & \statusopen \\
A finite-$q$ sweet spot (Prob.~\ref{prob:finite-q}) & \statusopen \\
\bottomrule
\end{tabular}
\end{center}

\noindent
$^{\dagger}$ modulo the real-scalar Hirschfeld--\.Zelazko input; see
Remark~\ref{rem:hz-status} and Problem~\ref{prob:real-hz}.

\subsection{On the numerical tags}

The searches in \texttt{verify/} maximise nonsmooth, nonconvex ratios over the
unit sphere of a finite-dimensional algebra.  Every reported value is therefore
a \emph{lower} bound for a supremum, refined until it stops moving; a
\statusnum\ tag means ``the search found nothing above this, from many
starts'', not ``this is proved''.  Three safeguards are built in.  The gauge
itself is also computed by maximisation, so an under-resolved gauge inflates a
ratio; every coarse search is followed by a refinement pass with a high-accuracy
gauge, because an early version of the code reported a spurious
$\kap(T_2(\RR))=0.317$ caused by exactly that.  For $T_2(\RR)$, where the
answer matters most, the gauge is computed exactly instead: every commutator in
$T_2(\RR)$ is a multiple of $E_{12}$, so $\gauge_1$ there is the square of a
support function, which by convex duality equals a one-dimensional minimisation
of a nuclear norm (Lemma~\ref{lem:dual-certificate}) --- no search at all.  And
where a numerical answer turned out to matter structurally, we went back and
proved it: the separation $\kap_{\gauge_1}(T_2(\RR))\ge\kap^*>\tfrac14$ began as
a number from the search and is now Proposition~\ref{prop:T2-constant}, with the
search retained only for the claim that the bound is attained.
See \texttt{verify/experiments/t2\_exact.py}.

\section{Open problems}
\label{sec:open}

Ordered by leverage: the first three would change the shape of the theory, the
rest would firm up statements already in use.

\begin{problem}[Completeness]\label{prob:completeness}
Decide whether $\PB^C_{\gauge_1}$ has equalizers.  By
Question~\ref{q:equalizer-exists} this is the representability of the
admissible-character-pair functor $F$, and by Theorem~\ref{thm:t2-cone} the
naive candidate $\RR1$ is eliminated.  Concretely: is there a PB-algebra $E$
carrying a universal admissible pair of characters?  A negative answer makes
$\PB$ incomplete --- the expected outcome --- and a positive one would produce
an interesting new object, namely the free PB-algebra on a pair of characters
whose separation is funded by its own commutators.
\end{problem}

\begin{problem}[Semisimplicity at the critical constant]\label{prob:semisimple}
Prove or refute Conjecture~\ref{conj:semisimple}: is every
$A\in\PB^{1/4}_{\gauge_1}$ semisimple?

This is the question opened up by Theorem~\ref{thm:constant-separates}, which
refutes the earlier notes' claim that no $\gauge_1$-constant excludes
$T_2(\RR)$.  A positive answer would be the strongest structural statement
available about the theory: at exactly one value of the constant, the norm gauge
would exclude the radical --- the thing the $C^*$-identity excludes by means of
an involution --- while keeping the functorial classical base that
$\gauge_\infty$ and every $\gauge_q$, $q\ge2$, destroy
(Warning~\ref{warn:vanishing}).  A counterexample would be nearly as useful: a
PB-algebra at $C=\tfrac14$ with nonzero radical would be the first object of the
theory that no $C^*$-algebra can imitate, and by
Warning~\ref{warn:cone-constant} it would also be the natural place to look for
a legal cone at the critical constant.

Concrete first steps: compute $C_{\gauge_1}(T_n(\RR))$ for $n\ge3$ and for the
other elementary radical-carrying algebras; then look for a general lower bound
on $\kap$ in terms of $\rad(A)$.
\end{problem}

\begin{problem}[Closure under quotients]\label{prob:quotients}
Is $\PB^C_{\gauge_1}$ closed under quotients by closed two-sided ideals?
Quotient maps are always legal (Proposition~\ref{prop:category}); at issue is
whether $A/I$ satisfies \eqref{ax:pb1} and \eqref{ax:pb2} with the same
constant.  The danger is that the quotient shrinks the budget faster than the
defect, since the budget is a supremum over test elements that the quotient may
collapse.  This single question controls two others: coequalizers, hence
cocompleteness (\S\ref{sec:categorical}), and whether the fibres $A_x$ of the
bundle picture are PB-algebras (\S\ref{sec:bundle}).  A counterexample would be
almost as valuable as a proof, since it would identify the extra hypothesis the
bundle picture needs.
\end{problem}

\begin{problem}[Coproducts]\label{prob:coproducts}
Does $\PB^C_{\gauge_1}$ have coproducts?  Warning~\ref{warn:not-free-product}
shows the free product is not the answer --- on commutative objects the
coproduct is $C(X\times Y,\RR)$, by Corollary~\ref{cor:comm-coproducts} --- so
the question is what replaces it in general.  A natural guess is the
``budget-completed'' free product: the free product quotiented by whatever is
needed to keep the images' gauges from growing.  Compute the first
noncommutative case, $M_2(\RR)\amalg M_2(\RR)$.
\end{problem}

\begin{problem}[The uniform $C^*$ bound]\label{prob:cstar-bound}
Prove or refute Conjecture~\ref{conj:cstar-uniform}:
$\norm{a}^2-\norm{a^2}\le\dist(a,\RR1)^2$ for every element of a real
$C^*$-algebra.  A proof puts every real $C^*$-algebra into $\PB^{1/4}$ as far as
\eqref{ax:pb1} is concerned, reducing membership in $\mathcal{I}$ to the
spectral-reality axiom alone, which is a condition on the centre.
\end{problem}

\begin{problem}[Real Stampfli]\label{prob:real-stampfli}
Prove Proposition~\ref{prop:real-stampfli} for $M_n(\RR)$, i.e.\
$\norm{\ad_a}=2\dist(a,\RR1)$ over the reals, or find the correct real
statement.  Stampfli's theorem is over $\CC$; the numerics of \texttt{verify/}
find equality for $n\le4$ but also show the identity \emph{fails} on
$T_n(\RR)$, so whatever is true is sensitive to more than the norm.
\end{problem}

\begin{problem}[Matrix stability]\label{prob:matrix-stability}
Is $\kap_{\gauge_1}(M_n(\RR))=\tfrac14$ for all $n$, and more generally is
$C_{\gauge_1}(M_n(A))$ bounded in $n$ for a fixed PB-algebra $A$?  Uniform
boundedness is the minimum requirement for any $K$-theoretic or Morita-theoretic
development, and without it the category cannot support the usual
noncommutative-geometry machinery.  The numerics confirm the first question for
$n\le4$.
\end{problem}

\begin{problem}[Quaternionic Stone--Weierstrass]\label{prob:quaternionic-sw}
Prove that \eqref{ax:sq} together with \eqref{ax:pb2} forces an isometric
unital embedding into an algebra of sections with $\RR$- and $\HH$-fibres, the
infinite-dimensional form of Proposition~\ref{prop:strict-survivors}.
\end{problem}

\begin{problem}[The finite-$q$ sweet spot]\label{prob:finite-q}
Is there a single finite $q$ whose class excludes $T_2(\RR)$ while remaining
closed under quotients?  Lemma~\ref{lem:gauge-basic}(e) shows the cost in
advance: for every $q\ge2$ the vanishing locus of $\gauge_q$ is strictly larger
than the centre, so Lemma~\ref{lem:centre-preserved} fails and the base is no
longer functorial.  Any positive answer should therefore come with a
replacement for \S\ref{sec:base}.
\end{problem}

\begin{problem}[Real Hirschfeld--\.Zelazko]\label{prob:real-hz}
Give a self-contained proof of the real-scalar input to
Theorem~\ref{thm:commutativity-forced}: a real Banach algebra with
$r=\norm{\cdot}$ is commutative.  See Remark~\ref{rem:hz-status} for what
exactly is needed; the obstacle is that $r=\norm{\cdot}$ on $A$ does not
transfer verbatim to $A_\CC$.
\end{problem}

\begin{problem}[An intrinsic budget]\label{prob:intrinsic}
Is there an element-intrinsic, involution-free replacement for $\gauge$ --- one
not defined as a supremum over the ambient algebra --- which agrees with
$\gauge_1$ on the examples?  By Remark~\ref{rem:ambient-referencing} this is
precisely what stands between $\PB$ and a limit theory: an intrinsic budget
would be inherited by subalgebras, the solution set $\{f=g\}$ would be an
object with a legal inclusion, and the $C^*$ proof of completeness would go
through verbatim.  We regard this as the single most consequential question on
the list, and note that its answer is plausibly negative --- in which case the
right theorem is a proof that no such replacement exists.
\end{problem}

\begin{problem}[Faithfulness of the bundle]\label{prob:faithful}
Answer Question~\ref{q:faithful-bundle}: is
$\norm{a}=\sup_{x\in X_A}\norm{\pi_x(a)}$?  Equivalently, does the central
$C(X_A,\RR)$-module structure satisfy the partition-of-unity estimate that
holds automatically for $C^*$-algebras?
\end{problem}


\bibliographystyle{plain}
\bibliography{refs}

\end{document}
